\documentclass{amsart}
\usepackage{amssymb,mathtools,array,longtable,booktabs}
\usepackage[hidelinks]{hyperref}
\setlength{\emergencystretch}{3em}
\raggedbottom
\hypersetup{pdftitle={Irrationality of the trilogarithm at reciprocals of integers of absolute value at least 27},pdfauthor={Ruiyi Zhang}}

\newtheorem{theorem}{Theorem}
\newtheorem{lemma}{Lemma}[section]
\newtheorem{proposition}[lemma]{Proposition}
\newtheorem{fact}{Fact}
\theoremstyle{remark}
\newtheorem{remark}{Remark}[section]
\numberwithin{equation}{section}
\numberwithin{table}{section}

\newcommand{\Li}{\operatorname{Li}}
\newcommand{\Z}{\mathbb{Z}}
\newcommand{\Res}{\operatorname*{Res}}
\newcommand{\dd}{\,\mathrm d}
\newcommand{\needlines}[1]{\par\dimen0=\pagegoal\advance\dimen0-\pagetotal
  \ifdim\dimen0<#1\baselineskip\newpage\fi}
\begin{document}

\title[Irrationality of $\operatorname{Li}_3(1/a)$ for $|a|\ge27$]{Irrationality of the trilogarithm at reciprocals of integers of absolute value at least 27}
\author{Ruiyi Zhang}
\address{Independent Researcher, Columbia, Maryland, USA}
\email{ryzfsu@gmail.com}
\subjclass[2020]{Primary 11J72; Secondary 33B30}
\keywords{trilogarithm, irrationality, determinant, logarithmic energy}

\begin{abstract}
We prove that $\operatorname{Li}_3(1/a)$ is irrational for every integer
$a\ge27$ and every integer $a\le-27$. One determinant construction is used
on both sides. The positive side takes $m=12Q$ and $g=4Q$; the negative side
takes $m=16Q$ and $g=5Q$. We prove a conservative coefficientwise
integerizer, nonvanishing along prime-power towers, and a
logarithmic-energy estimate. On the negative side the estimate rests on a
signed residue identity along a bent contour. Exact rational and directed
interval certificates give four strict endpoint margins. Slope comparisons
carry them along both rays.
\end{abstract}

\maketitle

\section{Introduction}

For $|z|<1$, define
\[
\operatorname{Li}_3(z)=\sum_{n\ge1}\frac{z^n}{n^3}.
\]

By Marcovecchio \cite[p.~229]{Marcovecchio}, Hata's method
\cite[Corollary~2.2]{Hata} gives the linear independence of
\[
1,\ \operatorname{Li}_1(1/a),\ \operatorname{Li}_2(1/a),\
\operatorname{Li}_3(1/a)
\]
for $|a|\ge1038$. Panzone gave a short proof that
$\operatorname{Li}_3(1/d)$ is irrational for every integer $d\ge1173$
\cite[Theorem~2, p.~108]{Panzone}. Marcovecchio proved the four-number
independence for $a\ge486$ and $a\le-471$ \cite[Theorem~7.1,
pp.~228--229]{Marcovecchio}. In these ranges $\operatorname{Li}_3(1/a)$ is
therefore irrational.

This paper continues the construction in \cite{Zh,ZhB}. Similar determinant constructions appear
in \cite{Fauzan,Lai}. For the dilogarithm at rational points see also
\cite{Prellberg,Liu,TanZhang}.

Our earlier preprint \cite{ZhB} treated $a\ge67$ and $a\le-70$. It used
$g=m/2$ and $m=4Q$. We cite its routine lemmas where they carry over, and we
list the changes each time. The new arithmetic rates, the seed towers, the
endpoint certificates and the signed negative contour are proved here. We
prove the following statement.

\begin{theorem}\label{thm:main}
For every integer $a\ge27$ and every integer $a\le-27$,
$\operatorname{Li}_3(1/a)$ is irrational.
\end{theorem}

The proof has three parts. Sections~\ref{sec:determinant}--\ref{sec:arithmetic}
give the determinant and its integerization. Section~\ref{sec:nonvanishing}
proves nonvanishing. Sections~\ref{sec:positive} and \ref{sec:negative}
give the analytic estimates for the two signs.
Sections~\ref{sec:endpoints}--\ref{sec:proof} assemble the four endpoint
certificates and the rays.

In Remark~\ref{rem:scope}, let $G_{m,a}$ be the determinant polynomial of
Section~\ref{sec:determinant}. For $G_{m,a}\not\equiv0$, let $\Delta_m$ be
the least positive rational number with $\Delta_mG_{m,a}\in\mathbb Z[T]$,
and put $\tau_m=m^{-2}\log\Delta_m$. When
$G_{m,a}(\operatorname{Li}_3(1/a))\ne0$, put
$\alpha_m=-m^{-2}\log|G_{m,a}(\operatorname{Li}_3(1/a))|$ and
$B_m=\tau_m-\alpha_m$. The criterion applies when these nonzero evaluations
contain a subsequence $m_k$ on which $B_{m_k}$ stays below a negative
constant. It also needs, for every fixed rational $\xi$,
$G_{m_k,a}(\xi)\ne0$ for all sufficiently large $k$. Finite values of
$B_m$ are only indicative. The ratio is $\rho=g/m$. The
staggered filter replaces $\binom{r-1}{g}^3$ by a product of three binomial
coefficients with different lower indices. This variant construction is not
used in this paper.

\begin{remark}[Scope of the present construction]\label{rem:scope}
We do not attempt to optimize the threshold in this paper. We expect that
the threshold could be lowered, but we make no lower-threshold claim here.

Several choices were made for safety and are not optimal. The
integerization constants are conservative. The filter ratio $\rho$, the
index windows and the negative-side contour tilt were not optimized jointly.
Our expectation is that tuning them could lower the threshold by a few more
units. The tuning would keep the same single-block determinant and the same
three proof steps:
integerization, nonvanishing and analytic decay.

The data suggest, however, that this algorithm is unlikely to reach
thresholds below 20. We computed exact denominators and determinants of the
same construction for $m\le128$. We used $g=5m/16$ and $g=\lceil m/3\rceil$.
We recorded the difference
$B_m=\tau_m-\alpha_m$ between the arithmetic growth and the analytic decay
per $m^2$:
\begin{itemize}
\item $a=+20$. $B_m$ crosses zero. For $g=\lceil m/3\rceil$ and $g=5m/16$
it equals $-0.0728$ and $-0.0494$ at $m=64$, and $+0.007$ and $+0.017$ at
$m=112$.
\item $a=+21,+22$. $B_m$ is still negative at $m=112$: about $-0.05$ to
$-0.06$ for $a=21$, and about $-0.10$ to $-0.11$ for $a=22$.
\item $a=-20$. For $g=5m/16$, $B_m$ stays negative but rises fast: about
$-0.21$ at $m=64$, $-0.12$ at $m=112$ and $-0.08$ at $m=128$.
\item Staggered filter $\binom{r-1}{m/8}\binom{r-1}{3m/16}\binom{r-1}{m/4}$.
It gives $B_m>0$ at $a=\pm20$ for $m=64,80,96,112$.
\end{itemize}
These finite computations put the apparent crossover of the method near
$|a|=21$--$23$ on the positive side. On the negative side it lies close to
20. We therefore expect that going below $|a|=20$ requires a structurally
different auxiliary construction. Further parameter tuning should not
suffice. All statements in this remark are numerical observations at finite
$m$. They are not part of the proof.
\end{remark}

\section{Determinant construction}\label{sec:determinant}

Fix an integer $a$ with $|a|\ge2$. We work in one of the two towers
below. In each, $m$ is a positive multiple of $4$ and $g<m$. Put
\[
 \rho=\frac gm,\qquad I=\{0,\ldots,m-1\},\qquad
 J=\{m,\ldots,2m-1\}.
\]
The positive tower is
\[
 (m,g)=(12Q,4Q),\qquad \rho=\frac13,
\]
and the negative tower is
\[
 (m,g)=(16Q,5Q),\qquad \rho=\frac5{16}.
\]
Here and below $Q$ is a positive integer. The nonvanishing section
specializes it to $Q=p^k$.
Define
\[
 w(r)=\binom{r-1}{g}^{3},\qquad
 F_{ij}(r)=\frac{w(r)}{(r+i)(r+j)},\qquad K_{ij}(r)=F_{ij}''(r),
\]
and
\[
 M_{ij}(a)=\sum_{r\ge1}a^{-r}K_{ij}(r),\qquad i\in I,\ j\in J.
\]
Polynomial division at the two poles gives
\[
 F_{ij}(r)=P_{ij}(r)+\frac{c_i}{r+i}+\frac{c_j}{r+j},\quad
 c_i=\frac{w(-i)}{j-i},\quad c_j=-\frac{w(-j)}{j-i}.
\]
At a pole of the partial fractions we take the finite part. Put
$b=a-1$, and define
\[
 \begin{aligned}
 E_{ij}(a)&=\sum_{r\ge1}a^{-r}P_{ij}''(r)
 -2c_i a^i\sum_{n=1}^{i}\frac{a^{-n}}{n^3}
 -2c_j a^j\sum_{n=1}^{j}\frac{a^{-n}}{n^3},\\
 C_{ij}(a)&=2c_i a^i+2c_j a^j,\qquad
 G_{m,a}(T)=\det(E(a)+TC(a)).
 \end{aligned}
\]
Empty sums are zero.

\begin{lemma}\label{lem:decomposition}
For every integer $a$ with $|a|\ge2$,
\[
 M(a)=E(a)+\operatorname{Li}_3(1/a)C(a).
\]
The polynomial $G_{m,a}$ has degree at most $m$.
\end{lemma}

\begin{proof}
Differentiating the partial fractions gives
\[
 F_{ij}''=P_{ij}''+\frac{2c_i}{(r+i)^3}
                    +\frac{2c_j}{(r+j)^3}.
\]
Also
\[
 \sum_{r\ge1}\frac{a^{-r}}{(r+i)^3}
 =a^i\left(\operatorname{Li}_3(1/a)
              -\sum_{n=1}^i\frac{a^{-n}}{n^3}\right).
\]
All series converge absolutely because $|a|>1$. Substitution gives the
matrix identity. Each matrix entry is affine in $T$. This proves the
degree bound. At a pole we differentiate the numerator before taking the
finite part. With $\delta=j-i$ this gives
\[
 P_{ij}''(-i)=\frac{w'''(-i)}{3\delta}
 -\frac{w''(-i)}{\delta^2}+\frac{2w'(-i)}{\delta^3}
 +\frac{2(w(-j)-w(-i))}{\delta^4},
\]
 and, with the same $\delta=j-i$,
\[
 P_{ij}''(-j)=-\frac{w'''(-j)}{3\delta}
 -\frac{w''(-j)}{\delta^2}-\frac{2w'(-j)}{\delta^3}
 +\frac{2(w(-i)-w(-j))}{\delta^4}.
\]
These are finite parts at the poles of the partial fractions. They are the
values used below.
\qedhere
\end{proof}

\section{The determinant criterion}\label{sec:criterion}

The following fact is the elementary limsup variant of the determinant
criterion in \cite[Sec.~3]{Zh}. We prove the exact form used here.
We use the extended-real convention $\log 0=-\infty$ in its limsup bounds.

\begin{fact}[Determinant criterion]\label{fact:criterion}
Let $\theta$ be real. For a fixed base $a$, let $G_{m,a}(T)$ have degree at
most $m$, and let $m_k$ increase. Suppose positive rationals $\Delta_m(a)$
satisfy
\[
 \Delta_m(a)G_{m,a}(T)\in\mathbb Z[T],\qquad
 \limsup_k m_k^{-2}\log\Delta_{m_k}(a)\le\tau(a),
\]
and
\[
 \limsup_k m_k^{-2}\log|G_{m_k,a}(\theta)|\le-\alpha(a),
 \qquad \alpha(a)>\tau(a).
\]
If every fixed rational $\xi$ satisfies
$G_{m_k,a}(\xi)\ne0$ for all sufficiently large $k$, then $\theta$ is
irrational.
\end{fact}

\begin{proof}
If $\theta=u/v$ in lowest terms, then
\[
 Z_m=v^m\Delta_m(a)G_{m,a}(u/v)
   =\sum_{\ell=0}^m(\Delta_m\gamma_{m,\ell})u^\ell v^{m-\ell}
\]
is a nonzero integer whenever
$G_{m,a}(T)=\sum_\ell\gamma_{m,\ell}T^\ell$ is nonzero at $u/v$.
The factor $v^m$ is $\exp(O(m))$. So
$\limsup m_k^{-2}\log|Z_{m_k}|\le\tau(a)-\alpha(a)<0$, and hence
$|Z_{m_k}|\to0$. This is a contradiction. In the application,
$\Delta_m=D_m(a)$ from \eqref{eq:9}. The prime-power tower in
Section~\ref{sec:nonvanishing} supplies the nonvanishing hypothesis.
\qedhere
\end{proof}

\section{Coefficient integerization}\label{sec:arithmetic}

The arithmetic construction is coefficientwise. It uses no primitive-content
or exact-gcd assertion. Put
\[
 U=2m-1,\quad L=\operatorname{lcm}(1,\ldots,U),\quad
 N=\frac{m^2}{8}-\frac m2,\quad H=3g-3,
\]
and
\[
 H_*=\left\lfloor\frac{(H+1)^2}{4}\right\rfloor
     =\left\lfloor\frac{(3g-2)^2}{4}\right\rfloor .
\]
For a prime $p$, let $\ell_p=v_p(L)$. Let
$\kappa_n$ be $1$ when
$n\bmod p\ge p-(g\bmod p)$, and $0$ otherwise.

\subsection{Formal divisors}

We use the coefficientwise Newton argument of Lemma~4.1 of \cite{ZhB},
together with its divided-difference lemmas. Newton interpolation, two divided
differences, and two derivatives give
\begin{equation}
 L^{4m}b^{mH}G_{m,a}(T)\in\mathbb Z[a,T].                 \label{eq:1}
\end{equation}
We use the present values of $g$. The four copies of $L$ account for the two
divided differences and the two derivatives. The pole and harmonic terms
need no extra divisor. Their complete denominators include
$1/(j-i)$ and $1/((j-i)n^3)$, respectively. These divide $L^4$ because
$|j-i|,n\le U$. In the Newton basis we use the
two denominator identities
\begin{equation}
 \sum_{r\ge1}a^{-r}\binom r k
   =\frac{a^{-k}}{(1-a^{-1})^{k+1}}\ (k\ge1),\qquad
 \sum_{r\ge1}a^{-r}\binom{r-1}{k}
   =\frac{a^{-k-1}}{(1-a^{-1})^{k+1}}.                    \label{eq:newton-denominators}
\end{equation}
Let $V$ be an integer-valued polynomial on $\mathbb Z^k$. Suppose it has
degree at most $D$ in one variable $X$. Put
$L_D=\operatorname{lcm}(1,\ldots,D)$. Then
$L_D\,\partial_XV$ and $L_D\,(V(X,\cdot)-V(Y,\cdot))/(X-Y)$ are integer
valued.
It suffices to take $V=\binom Xn U$ with $n\le D$ and $U$ integer valued in
the other variables. Vandermonde's identity gives
$\binom Xn-\binom Yn=\sum_{k\ge1}\binom{X-Y}k\binom Y{n-k}$, and
$\binom{X-Y}k/(X-Y)=\binom{X-Y-1}{k-1}/k$. For the derivative,
$\frac{\dd}{\dd X}\binom Xn=\sum_{k=1}^n\frac{(-1)^{k-1}}k\binom X{n-k}$.
Each of the two divided differences in $P_{ij}''$ therefore costs one
factor $L_D$. So does each of the two derivatives. Here $D\le U$.
With $L_D\mid L$ this accounts for the four copies of $L$ in \eqref{eq:1}. Thus a degree-$d$ Newton expansion
introduces only denominators dividing the appropriate power of
$L_d=\operatorname{lcm}(1,\ldots,d)$.
The identity is formal in $a$ and remains valid at negative integer bases.

Write $d=(g!)^3$ and take a polynomial $P\in\mathbb Z[1/d][X]$. Its
rational generating function has two expansions
\begin{equation}
 \sum_{r\ge1}a^{-r}P(r)=-\sum_{k\ge0}a^kP(-k),                 \label{eq:formal-expansions}
\end{equation}
where the left side expands at infinity and the right side at zero.
In the determinant calculation the right-hand expansion is
used after multiplication by the factor $a^N$ below.
Thus the displayed identity is an identity of formal power series at the
required order.
For $W\in\mathbb Z[X]$ we also use two divided-difference facts
\begin{equation}
 \frac{W(X)-W(Y)}{X-Y},\qquad
 \frac{W'(X)-W'(Y)}{X-Y}\in\mathbb Z[X,Y].                    \label{eq:divided-differences}
\end{equation}
Take $W=d w$. The divided differences in the partial fractions then lie
in $\mathbb Z[1/d][X,Y]$. In particular every coefficient of
$P_{ij}(r)$ is a polynomial in $(i,j)$ over $\mathbb Z[1/d]$.

Translate $T=X+\ell(a)$, where
$\ell(a)=\sum_{s\ge1}a^s/s^3$. Away from the node sets, the coefficient of
$a^k$ in $E+(X+\ell(a))C$ is $-F_{ij}''(-k)$. Away from the node sets it
has rank at most three. After two differentiations it is a sum of the
three Cauchy features $(i-k)^{-1},(i-k)^{-2},(i-k)^{-3}$. At a node the
finite part adds one exceptional feature. Suppose
$k\in I$ is the node, and put $\delta=j-k$ for $j\in J$. The
finite-part formula and the harmonic correction give
\begin{equation}
\begin{aligned}
 [a^k]\bigl(E+(X+\ell(a))C\bigr)_{k j}
  ={}\!&-\frac{w'''(-k)}{3\delta}+\frac{w''(-k)}{\delta^2}
      -\frac{2w'(-k)}{\delta^3}\\
 &\quad+\frac{2w(-k)}{\delta^4}
      +\frac{2Xw(-k)}{\delta}.
\end{aligned}
\label{eq:node-coefficient}
\end{equation}
The transposed formula holds when the node is in $J$. Thus a node row is a
linear combination of the four Cauchy row features
$\delta^{-1},\delta^{-2},\delta^{-3},\delta^{-4}$. The $X$ term is absorbed
in the coefficient of $\delta^{-1}$. A non-node row uses only the first
three. This proves the rank bound four, including the finite part. It is the
calculation above. Thus at most four columns per power
$a^k$ can contribute.
The order is therefore at least
\[
 \sum_{j=0}^{m-1}\left\lfloor\frac j4\right\rfloor
 =4\sum_{r=0}^{m/4-1}r=N.
\]
Consequently coefficientwise multilinearity gives
\[
 a^N\mid\det(E+(X+\ell(a))C).
\]
Since $\ell(0)=0$, translating back preserves this divisibility.

At $b=0$ we use the following degree count. The polynomial $dw$ is monic of
degree $3g$. After polynomial division and two derivatives, $P_{ij}''(X)$
therefore has total degree at most $3g-4$ in $(X,i,j)$. Equivalently,
\[
 P_{ij}''(r)=d^{-1}\sum_{t=0}^{3g-4}p_t(i,j)r^t,
 \qquad \deg_{i,j}p_t\le 3g-4-t.
\]
The standard identity for $\sum_{r\ge1}a^{-r}r^t$ has denominator
$b^{t+1}$. Hence the coefficient of $b^{-s}$ receives terms only from
$t\ge s-1$. Its total degree in $(i,j)$ is at most
$3g-4-(s-1)=H-s$. The finite harmonic terms and $C$ are regular at
$b=0$. Taking forward differences of orders $u$ and $v$ lowers the pole
order to $(H-u-v)_+$. Suppose $k$ entries of a determinant term have
a positive pole bound. The distinct row and column orders then cost at
least $k(k-1)$. Hence the total pole order is at most
\[
 k(H+1-k)\le H_*.
\]
This is the forward-difference argument of \cite[Lemma~4.8]{ZhB}. There $H_m$ is
replaced by
$H_* =\max_{0\le k\le H+1}k(H+1-k)$. The harmonic terms and $C$ have no
pole at $b=0$. The row and column operations are unimodular. Thus
\begin{equation}
 b^{mH-H_*}\mid L^{4m}b^{mH}G_{m,a}(T).                   \label{eq:2}
\end{equation}
The polynomials $a$ and $a-1$ are monic and coprime. So
\eqref{eq:1}--\eqref{eq:2} give the integral formal quotient
\[
G_{m,a}(T)=\frac{a^N R_m(a,T)}{L^{4m}(a-1)^{H_*}},
\qquad R_m\in\mathbb Z[a,T].
\]

\subsection{Local prime weights}

The following weights are valid coefficientwise for the pencil $E+TC$,
including finite parts. The branches are applied in the displayed order.
The low-prime rank branch assumes $p\nmid a(a-1)$. If
$p\mid a(a-1)$ we use the endpoint weight $e_p^{\rm end}$ defined just before
\eqref{eq:9} instead. Primes $p>U$ have weight
zero. If $p<5$ or $p^2\le U+g$, use $4m\ell_p$. Now let $p\le g$, $p\ge5$,
and $p^2>U+g$. The entry demand is
\begin{equation}
 \mathfrak b_{ij}=3-\kappa_i\kappa_j+
             (1-\kappa_i){\bf1}_{i\equiv j\pmod p}.       \label{eq:3}
\end{equation}
The row and column weights $2-\kappa_n$ dominate this demand.

The saving rests on the following coefficientwise lemma. Its proof
includes finite parts and the $T$ coefficient.
% Insert for the low-prime rank lemma.
% This file is intentionally separate from every frozen paper directory.

\subsection{The low-prime rank lemma}
\label{sec:insert-low-prime-rank}

This subsection gives the coefficientwise local argument for the low-prime
branch. It applies to the pencil $E+TC$, including finite parts at poles and
the coefficient of $T$.
Throughout this subsection
\[
 I=\{0,\ldots,m-1\},\qquad J=\{m,\ldots,2m-1\},\qquad U=2m-1,
\]
\[
 w(X)=\frac{\mathcal N(X)}d,\qquad
 \mathcal N(X)=\prod_{u=1}^{g}(X-u)^3,\qquad d=(g!)^3.
\]
Thus $w(r)=\binom{r-1}{g}^{3}$ at every non-negative integer $r$.

\begin{lemma}[Low-prime rank]
\label{lem:residue-copy-rank-insert}
Let $p\geq5$ satisfy
\[
 p^2>U+g,
 \qquad p\nmid a(a-1),
 \qquad q=\left\lfloor\frac gp\right\rfloor\geq1,
\]
where $a$ is an integer with $|a|>1$.  Write
\[
 g=qp+h,\qquad 0\leq h<p,\qquad G=p-h.
\]
Call a node $n$ good if $n\bmod p<G$ and bad if
$n\bmod p\geq G$.  Put
\[
 \kappa_n=\begin{cases}0,&n\text{ good},\\1,&n\text{ bad},\end{cases}
 \qquad
 \omega_i=2-\kappa_i,\qquad \omega'_j=2-\kappa_j.
\]
For $n\geq0$ define
\[
 R_n(X)=\frac{\mathcal N(X)-\mathcal N(-n)}{X+n},
 \qquad
 H_n(a)=\sum_{s=1}^{n}\frac{a^{n-s}}{s^3},
\]
and
\begin{equation}
 b_n(T)=\sum_{r\geq1}a^{-r}\left(\frac{R_n}{d}\right)''(r)
       +2w(-n)\bigl(a^nT-H_n(a)\bigr).
 \label{eq:insert-bn}
\end{equation}
Let the normalized pencil entries be
\begin{equation}
 B_{ij}(T)=\frac{b_i(T)-b_j(T)}{j-i},
 \qquad i\in I,\ j\in J.
 \label{eq:insert-Bij}
\end{equation}
By the second divided-difference identity and the finite harmonic terms,
these normalized entries are exactly the entries of the determinant pencil:
\[
 B_{ij}(T)=E_{ij}(a)+TC_{ij}(a).
\]
Then the exact row and column scaling is
\begin{equation}
 \widehat B_{ij}(T)=p^{\omega_i+\omega'_j}B_{ij}(T)
 =p^{(2-\kappa_i)+(2-\kappa_j)}B_{ij}(T).
 \label{eq:insert-scaled-B}
\end{equation}
The matrix $\widehat B$ has entries in $\mathbb Z_{(p)}[T]$, and, if
$N_I$ and $N_J$ denote the numbers of good nodes in $I$ and $J$, respectively,
\begin{equation}
 \operatorname{rank}_{\mathbb F_p(T)}\overline{\widehat B}
 \leq \min(N_I,N_J)+\min(G,h).
 \label{eq:insert-rank-bound}
\end{equation}
Here the bar denotes reduction in the Gauss valuation ring of
$\mathbb Q_p(T)$.  Consequently, with
\[
 W_{\rm nat}=\sum_{i\in I}\omega_i+\sum_{j\in J}\omega'_j
             =2m+N_I+N_J,
\]
and put
\[
 \varrho_*:=\min(N_I,N_J)+\min(G,h).
\]
every coefficient of the determinant satisfies
\begin{equation}
 v_p\!\left([T^k]\det B(T)\right)
 \geq -W_{\rm nat}+m-\varrho_*
 \qquad(0\leq k\leq m).
 \label{eq:insert-rank-valuation}
\end{equation}

There is also a separate entrywise alternative.  Its row and column base
weights are $1$ and $2$, respectively.  For each good residue $s<G$, add one
extra power of $p$ to the side having fewer copies of that residue.  If
\[
 c_I(s)=\#\{i\in I:i\equiv s\pmod p\},\qquad
 c_J(s)=\#\{j\in J:j\equiv s\pmod p\},
\]
the resulting total exponent is
\begin{equation}
 W_{\rm alt}=3m+\sum_{0\leq s<G}\min\{c_I(s),c_J(s)\}.
 \label{eq:insert-W-alt}
\end{equation}
Either alternative alone supplies an integerizer; their savings are not
combined.
\end{lemma}

\begin{proof}
The proof has three steps: summation, local order, and rank.

\paragraph{\emph{Newton summation modulo $p$.}}
We first need an elementary summation fact.  Suppose
\[
 V(X)=\sum_{\ell=0}^{D}v_\ell\binom X\ell,
 \qquad v_\ell\in\mathbb Z_{(p)}.
\]
The ordinary series $\sum_{r\geq1}a^{-r}V(r)$ converges because $|a|>1$.
We claim that it lies in $\mathbb Z_{(p)}$.  Write $z=a^{-1}$.  Then
\begin{equation}
 \sum_{r\geq1}z^r\binom r0=\frac z{1-z},
 \qquad
 \sum_{r\geq1}z^r\binom r\ell
 =\frac{z^\ell}{(1-z)^{\ell+1}}\quad(\ell\geq1).
 \label{eq:insert-newton-sums}
\end{equation}
The only denominators are powers of $1-a^{-1}$.  This is a $p$-unit because
$p\nmid a(a-1)$.  Formula \eqref{eq:insert-newton-sums} also shows that
reduction modulo $p$ may be done term by term in the finite Newton
expansion.  We use no $p$-adic convergence of the infinite series.

We also need a periodic form.  Let $V$ be a polynomial with
$p$-integral Newton coefficients.  Suppose that
$V(r+p^2)\equiv V(r)\pmod p$ for every integer $r$.  Then
\begin{equation}
 \sum_{r\geq1}a^{-r}V(r)
 \equiv
 \frac{\displaystyle\sum_{r=1}^{p^2}a^{-r}V(r)}
      {1-a^{-p^2}}\pmod p.
 \label{eq:insert-periodic-sum}
\end{equation}
To see this, write $S=\sum_{r\geq1}a^{-r}V(r)$ and
$V(r+p^2)=V(r)+pV_1(r)$.  The polynomial $V_1$ is integer valued at
$p$-adic integers.  So it has $p$-integral Newton coefficients.  By the
first fact its sum is $p$-integral.  Splitting off the first $p^2$ terms gives
$S=\sum_{r=1}^{p^2}a^{-r}V(r)+a^{-p^2}(S+pS_1)$ with $S_1$ $p$-integral.
The denominator $1-a^{-p^2}$ is a unit since
$a^{-p^2}\equiv a^{-1}\not\equiv1\pmod p$.

\paragraph{\emph{Local order.}}
For a residue class $t\pmod p$, let
\[
 \mathfrak m_t=(p,X-t)\subset\mathbb Z_{(p)}[X]
\]
and let $\operatorname{ord}_t(f)$ be the largest $e$ for which
$f\in\mathfrak m_t^e$.  The local ring at $\mathfrak m_t$ is regular of
dimension two.  Its associated graded ring is
\begin{equation}
 \operatorname{gr}_{\mathfrak m_t}\mathbb Z_{(p)}[X]
 \simeq\mathbb F_p[P,Y],
 \qquad P=\operatorname{in}(p),\quad Y=\operatorname{in}(X-t).
 \label{eq:insert-associated-graded}
\end{equation}
This graded ring is an integral domain.  Three rules follow at once:
\begin{enumerate}
\item exact division by $X-t$ lowers local order by one.  Division by
  $X-u$ with $u\not\equiv t\pmod p$ is division by a unit.  It does not
  lower local order;
\item differentiation lowers local order by at most one;
\item if $x\equiv t\pmod p$, then evaluation gives
  $v_p(f(x))\geq\operatorname{ord}_t(f)$.
\end{enumerate}
The second rule follows by differentiating monomials
$p^r(X-t)^s$.  The third follows by substituting
$X-t=p\,u$ with $u\in\mathbb Z_p$.

Since $g<p^2$, Legendre's formula gives
\begin{equation}
 v_p(d)=3v_p(g!)=3q.
 \label{eq:insert-d-valuation}
\end{equation}
Every residue class occurs $q$ times among $1,\ldots,g$.  The classes
$1,\ldots,h$ occur once more.  Therefore
\begin{equation}
 \operatorname{ord}_t\mathcal N\geq3q,
 \qquad
 \operatorname{ord}_t\mathcal N\geq3q+3
 \quad\text{if }t\equiv1,\ldots,h\pmod p.
 \label{eq:insert-N-orders}
\end{equation}
If $n$ is bad, then $-n\equiv1,\ldots,h\pmod p$.  So
$\mathcal N(-n)$ has valuation at least $3q+3$.  If $n$ is good, its
valuation is at least $3q$.

\paragraph{\emph{Scaled node sums.}}
Take a residue $t\not\equiv-n\pmod p$.  There the denominator $X+n$ is a
unit, and \eqref{eq:insert-N-orders} gives
\[
 \operatorname{ord}_tR_n\geq3q.
\]
At $t\equiv-n\pmod p$ division lowers order by one.  The numerator then
has stronger bounds.  For bad $n$ both terms in the numerator have order at
least $3q+3$.  Hence $\operatorname{ord}_tR_n\geq3q+2$.  For good $n$ the
bound is $\operatorname{ord}_tR_n\geq3q-1$.  We take two derivatives and
divide by $d$.  The local-order rules then give
\begin{equation}
 p^2\left(\frac{R_n}{d}\right)''(r)\in\mathbb Z_p
 \quad(n\text{ bad}),
 \qquad
 p^3\left(\frac{R_n}{d}\right)''(r)\in\mathbb Z_p
 \quad(n\text{ good})
 \label{eq:insert-Rn-values}
\end{equation}
for every integer $r$.  By the finite-difference characterization of
$\operatorname{Int}(\mathbb Z,\mathbb Z_p)$, the scaled polynomials have
$p$-integral Newton coefficients.  So the Newton summation fact applies to
their complete sums.

Now take the finite-part term.  By hypothesis $n<U+1<p^2$.  So
$v_p(s)\leq1$ for $1\leq s\leq n$, and
$v_p(H_n(a))\geq-3$.  Also
\begin{equation}
 v_p(w(-n))\geq3\quad(n\text{ bad}),
 \qquad v_p(w(-n))\geq0\quad(n\text{ good}).
 \label{eq:insert-harmonic-bounds}
\end{equation}
So the $T$ and harmonic terms in \eqref{eq:insert-bn} have the same
scaled integrality as in \eqref{eq:insert-Rn-values}.  We have proved
\begin{equation}
 p^2b_n(T)\in\mathbb Z_{(p)}[T]\quad(n\text{ bad}),
 \qquad
 p^3b_n(T)\in\mathbb Z_{(p)}[T]\quad(n\text{ good}).
 \label{eq:insert-bn-integrality}
\end{equation}

\paragraph{\emph{The one-pole reduction.}}
Suppose that $i$ is bad and $j$ is good.  Their residues are distinct.  So
$j-i$ is a $p$-unit.  The first congruence in
\eqref{eq:insert-bn-integrality} gives
$p^3b_i(T)\equiv0\pmod p$.  The coefficient of $T$ in $p^3b_j(T)$ is
also zero modulo $p$, because $p^3w(-j)a^j\equiv0\pmod p$.  Define
\begin{equation}
 \mathfrak d_j(a)=p^3b_j(0)\bmod p.
 \label{eq:insert-Dj}
\end{equation}
It depends on neither $i$ nor $T$.  Since $p^3b_i(T)\equiv0$,
\eqref{eq:insert-Bij} gives $p^3B_{ij}(T)\equiv-\mathfrak d_j(a)/(j-i)$.  That is,
\begin{equation}
 p^3B_{ij}(T)\equiv \frac{\mathfrak d_j(a)}{i-j}\pmod p
 \qquad(i\text{ bad},\ j\text{ good}).
 \label{eq:insert-one-pole}
\end{equation}
The transpose formula holds for a good row and a bad column.

\paragraph{\emph{The bad--bad reduction.}}
Let $i$ and $j$ both be bad.  Put
\begin{equation}
 \delta=j-i,\qquad
 L_{ij}(X)=
 \frac{\mathcal N(-i)(X+j)-\mathcal N(-j)(X+i)}{j-i}.
 \label{eq:insert-Lij}
\end{equation}
Then the exact polynomial identity
\begin{equation}
 \mathcal N(X)-L_{ij}(X)=(X+i)(X+j)Q_{ij}(X)
 \label{eq:insert-Qij}
\end{equation}
defines $Q_{ij}\in\mathbb Z_{(p)}[X]$.  Since $0<|\delta|<2m<p^2$, we have
$v_p(\delta)\leq1$.  The two endpoint values of $\mathcal N$ have
valuation at least $3q+3$.  So every coefficient of $L_{ij}$ satisfies
\begin{equation}
 v_p\bigl([X^\ell]L_{ij}\bigr)\geq3q+2.
 \label{eq:insert-Lij-valuation}
\end{equation}

At an evaluation residue the vanishing factors fall into three cases.
Write $t\equiv r\pmod p$.  First suppose that the opposite node residue
$s\equiv-r\pmod p$ is bad.  Then
$\operatorname{ord}_t\mathcal N\geq3q+3$.
\begin{enumerate}
\item If both $X+i$ and $X+j$ vanish at $t$, then $i\equiv j\pmod p$ and
  $v_p(\delta)=1$.  Use the coordinate $X=t+Y$.  The constant coefficient
 of $L_{ij}$ has valuation at least $3q+3$.  Here the extra factor $t+i$
  or $t+j$ supplies one $p$ before division by $\delta$.  The coefficient
  of $Y$ has valuation at least $3q+2$.  Hence
  $\operatorname{ord}_tL_{ij}\geq3q+3$.
  Dividing by two vanishing factors gives
  $\operatorname{ord}_tQ_{ij}\geq3q+1$.
\item If exactly one factor vanishes, the other is a unit.  From
  \eqref{eq:insert-Lij-valuation} and the order $3q+3$ of $\mathcal N$,
  division by the one vanishing factor gives
  $\operatorname{ord}_tQ_{ij}\geq3q+1$.
\item If neither factor vanishes, both are units.  Equations
  \eqref{eq:insert-N-orders} and \eqref{eq:insert-Lij-valuation} give
  $\operatorname{ord}_tQ_{ij}\geq3q+2$.
\end{enumerate}
Divide by $d$ and differentiate twice.  Then multiplication by $p^2$ makes
each of the three cases divisible by $p$.  Thus
\begin{equation}
 p^2\left(\frac{Q_{ij}}d\right)''(r)\equiv0\pmod p
 \quad\text{when }-r\text{ is a bad residue}.               
 \label{eq:insert-Q-bad-residue}
\end{equation}

Now let $s\equiv-r\pmod p$ be good.  Both $X+i$ and $X+j$ are units in the
local ring at $t=-s$.  Set
$u(X)=((X+i)(X+j))^{-1}$.  The derivatives $u,u',u''$ are integral there.
All coefficients of $L_{ij}/d$ are divisible by $p^2$.  So the $L_{ij}$ term
vanishes after the present reduction.  Consider the terms of $(wu)''$
containing $w$ and $w'$.  By the local-order bounds for $w,w',w''$, they also
vanish after multiplication by $p^2$.  Consequently
\begin{equation}
 p^2\left(\frac{Q_{ij}}d\right)''(r)
 \equiv\frac{p^2w''(r)}{(r+i)(r+j)}\pmod p
 \quad\text{when }-r\text{ is good}.                         
 \label{eq:insert-Q-good-residue}
\end{equation}

The identity \eqref{eq:insert-Qij} is exactly the divided-difference identity
\[
 \frac{R_i(X)-R_j(X)}{j-i}=Q_{ij}(X).
\]
Therefore \eqref{eq:insert-Bij} takes the following form:
\begin{multline}
 B_{ij}(T)=\sum_{r\geq1}a^{-r}\left(\frac{Q_{ij}}d\right)''(r)\\
 +\frac{2w(-i)(a^iT-H_i(a))-2w(-j)(a^jT-H_j(a))}{j-i}.
 \label{eq:insert-B-Q}
\end{multline}
For bad $i,j$ consider the second term in \eqref{eq:insert-B-Q}.  We have
$v_p(w(-i)),v_p(w(-j))\geq3$, $v_p(H_i),v_p(H_j)\geq-3$ and
$v_p(j-i)\leq1$.  So its $T$ part has valuation at least $2$.  Its
harmonic part has valuation at least $-1$.  After multiplication by $p^2$
both vanish modulo $p$.

It remains to sum the good-residue contribution.  Define
\[
 \varpi(r)=p^2w''(r)\pmod p.
\]
The polynomial $p^2w''$ has $p$-integral Newton coefficients.  Its values
have period $p^2$ modulo $p$.  To see this, apply Taylor expansion and the
local-order rule.  For every $\ell\geq1$ they give
\begin{equation}
 v_p\!\left(\frac{p^{2\ell+2}}{\ell!}
              w^{(\ell+2)}(r)\right)
 \geq \ell-v_p(\ell!)\geq1,
 \label{eq:insert-z-period}
\end{equation}
where the last inequality uses $p\geq5$.  Now take $V=p^2(Q_{ij}/d)''$.
By the case analysis above its values are $p$-integral.  So it has
$p$-integral Newton coefficients.  By
\eqref{eq:insert-Q-bad-residue} and \eqref{eq:insert-Q-good-residue},
$V(r)\equiv \varpi(r)/((r+i)(r+j))$ when $-r$ is good and $V(r)\equiv0$
otherwise.  The factor $1/((r+i)(r+j))$ modulo $p$ depends only on
$r\bmod p$.  So $V(r+p^2)\equiv V(r)\pmod p$.  Hence
\eqref{eq:insert-periodic-sum} applies to $V$.  For a good residue $s<G$, set
\begin{equation}
 \mathfrak c_s(a)=\frac{\displaystyle
       \sum_{\substack{1\leq r\leq p^2\\-r\equiv s\ (p)}}a^{-r}\varpi(r)}
       {1-a^{-p^2}}\pmod p.
 \label{eq:insert-Cs}
\end{equation}
Grouping the terms in \eqref{eq:insert-B-Q} by $s\equiv-r\pmod p$ gives
\begin{equation}
 p^2B_{ij}(T)
 \equiv \sum_{0\leq s<G}
       \frac{\mathfrak c_s(a)}{(i-s)(j-s)}\pmod p
 \qquad(i,j\text{ bad}).
 \label{eq:insert-bad-bad}
\end{equation}
All displayed denominators are $p$-units, since $i,j$ are bad and $s$ is
good.  The expression is independent of $T$.

\paragraph{\emph{Rank count and valuation transfer.}}
Take a bad row $i$ and a good column $j$.  Then the reduction
\eqref{eq:insert-one-pole} is a Cauchy feature with the good residue $s\equiv j\pmod p$.
For a bad row and a bad column, \eqref{eq:insert-bad-bad} is a sum of the
same features $1/(i-s)$, $s<G$.  Now take the block of bad rows, with all
columns.  It factors through the $G$ good residue features.  So its rank is
at most $G$.  Bad rows with the same residue modulo $p$ are also identical in
every column.  For good columns this follows from
\eqref{eq:insert-one-pole}.  For bad columns it follows from
\eqref{eq:insert-bad-bad}.  So the rank is also at most $h$, the number of
bad residue classes.  The bad-row block has rank at most $\min(G,h)$.
Adding the $N_I$ good rows gives
\begin{equation}
 \operatorname{rank}\overline{\widehat B}
 \leq N_I+\min(G,h).
 \label{eq:insert-rank-I}
\end{equation}
The same argument for the transpose gives the bound
$N_J+\min(G,h)$.  Taking the smaller of the two proves
\eqref{eq:insert-rank-bound}.

The factors in \eqref{eq:insert-scaled-B} are the row and column
weights.  Bad nodes receive one power and good nodes two powers.  Thus
\begin{equation}
 \det\widehat B=p^{W_{\rm nat}}\det B.
 \label{eq:insert-det-scaling}
\end{equation}
The Gauss valuation on $\mathbb Q_p(T)$ has residue field $\mathbb F_p(T)$.
Let a matrix have residue rank at most $r$.  Gaussian elimination over the
valuation ring shows that its determinant has valuation at least $m-r$.
Equivalently, every $T$-coefficient of the determinant has valuation at
least $m-r$.  We apply this to $\widehat B$.  With
\eqref{eq:insert-rank-bound} it proves \eqref{eq:insert-rank-valuation}.

Finally, the one-pole bounds give $p^3B_{ij}\in\mathbb Z_{(p)}[T]$ for
distinct residues.  A same-good-residue pair can lose one further power in
the divided difference.  It is covered by $p^4B_{ij}$.  The bad--bad and
mixed cancellations above show that $p^3B_{ij}$ is integral in all remaining
cases.  So one row power and two column powers cover every edge except the
same-good-residue edges.  For a good residue $s$ only those edges need one
extra power.  Put that extra power on the side with fewer copies.  This gives
\eqref{eq:insert-W-alt} and completes the proof.
\qedhere
\end{proof}


For the $p>g$ branch, write
\[
 \widetilde c_{i,j}=\frac{w(-i)}{j-i},\qquad
 \widetilde c_{j,i}=-\frac{w(-j)}{j-i}
\]
for the normalized pole coefficients in the $(i,j)$ entry. To avoid
overloading the low-prime symbol $\kappa_n$, put
\[
 \chi_n={\bf1}_{\{n\bmod p\ge p-g\}}.
\]
In the non-generic branch assume $p^2>U+g$ and $p\nmid a(a-1)$. The
remaining small primes use the first line of \eqref{eq:W-cons-piecewise}.
Since $p>g$, $v_p(g!)=0$, and
\[
 w(-n)=(-1)^{3g}\frac{\prod_{s=1}^{g}(n+s)^3}{(g!)^3}.
\]
The interval of factors has length $g<p$, so
\[
 v_p(w(-n))\ge3\chi_n,
 \qquad v_p(j-i)=\varepsilon_{ij}\le1.
\]
The divided-difference formulas for $P_{ij}$ have coefficients in
$\mathbb Z_{(p)}[r]$; hence $P_{ij}\in\mathbb Z_{(p)}[r]$. Also
$v_p(s)\le1$ for $1\le s\le U$, and therefore
\[
 v_p(H_n(a))\ge-3{\bf1}_{n\ge p},\qquad
 v_p(\widetilde c_{i,j})\ge3\chi_i-\varepsilon_{ij},\qquad
 v_p(\widetilde c_{j,i})\ge3\chi_j-\varepsilon_{ij}.
\]
The $C$ coefficient is a sum of the two pole terms, so
\[
 v_p(C_{ij})\ge3\min(\chi_i,\chi_j)-\varepsilon_{ij}.
\]
Thus the entry demand is
\[
\mathfrak b_{ij}^{>g}=\max\{0,\,3\mathbf1_{i\ge p}-3\chi_i+\varepsilon_{ij},\,
                    3\mathbf1_{j\ge p}-3\chi_j+\varepsilon_{ij}\}.
\]
The samples are $0\le n\le3g-4<m$, so a pole $n=j$ is never sampled. The
entrywise bounds above and the two divided-difference formulas give the
coefficientwise valuation transfer for this $p>g$ branch. The remaining
expansion is a finite polynomial calculation. The displayed edge demand
\eqref{eq:3} refers instead to the low-prime branch $p\le g$ above. So do
the row and column weights $2-\kappa_n$ and the coefficientwise determinant
bound in \eqref{eq:insert-rank-valuation}. They follow from the local-order
calculation in Lemma~\ref{lem:residue-copy-rank-insert}.

For $g<p<m$, put $u=p-g$. The potential calculation of
\cite[Lemma~4.5 and Lemma~4.6, case (4)]{ZhB} is reproduced here in its
needed form. There are five edge types: active--active, active--killed,
low--active, low--killed, and killed--killed. Their demands are bounded above,
respectively, by $4$, $3$, $4$, $0$, and $\varepsilon_{ij}$, where
$\varepsilon_{ij}=v_p(j-i)\in\{0,1\}$.
A node in $[0,u)$ is active and one in
$[u,p)$ is killed. A node $n\ge p$ is active if its residue is active.
Call a row or column low when its index is $<u$; nodes in $[u,p)$ are
killed. Let $R$ and $C$ count active rows and columns whose indices are
$\ge p$. Let
\[
 \Sigma_{\rm kill}=\sum_{z=p-g}^{p-1}\min(c_I(z),c_J(z)),
 \quad c_I(z)=\#\{i\in I:i\equiv z\pmod p\},
\]
where $c_J(z)=\#\{j\in J:j\equiv z\pmod p\}$. The two valid potentials are
\begin{equation}
 W_A=3(R+C)+u+\Sigma_{\rm kill},\qquad
 W_B=m+2(R+C)+u.                                         \label{eq:5}
\end{equation}
In $W_A$, active rows and columns with index at least $p$ receive three
units. The $u$ low active rows contribute the displayed $u$. The smaller
side of each killed residue receives one extra unit. In $W_B$, active rows
with index at least $p$ receive two units, low rows one, and killed rows zero.
Active columns receive three and killed columns one. The edge demands are
covered in the five cases active--active, active--killed, low--active,
low--killed, and killed--killed. For $m\le p\le U$, put
\[
 r_p=2m-p,\qquad h_p=(2m+g-2p)_+,
\]
and use the row and column weights
\[
 \omega_i^{\rm tail}={\bf1}_{i<r_p}\quad(i\in I),\qquad
 \omega_j^{{\rm tail}\prime}=3\,\mathbf{1}_{\{p\le j<2m,\ j-p<p-g\}}
 \quad(j\in J).
\]
There are $r_p$ same-residue edges $j=i+p$, covered by their row units.
The positive nonsame-residue demands have a good column and are covered by
its three units. The number of such columns is $r_p-h_p$. Thus the total
weight is
\begin{equation}
 W_{\rm tail}=\sum_{i\in I}\omega_i^{\rm tail}
 +\sum_{j\in J}\omega_j^{{\rm tail}\prime}
 =r_p+3(r_p-h_p)=4(2m-p)-3(2m+g-2p)_+.                    \label{eq:6}
\end{equation}

For the low-prime profile, let $x=p/m$ and write
$\nu=\lfloor g/p\rfloor$; thus
$\rho/(\nu+1)<x<\rho/\nu$. Put
\[
 \gamma=(\nu+1)x-\rho,\quad\eta=\rho-\nu x,\quad
 k_\ell=\lfloor\ell/x\rfloor,
\]
\[
 A_\ell=k_\ell\gamma+\min(\ell-k_\ell x,\gamma),\quad
 A=A_1,\quad B=A_2-A_1,
\]
and
\[
 D(x)=\int_0^\gamma\min(c_I(t),c_J(t))\,\dd t,
\]
where, explicitly,
\[
 c_I(t)=\#\{k\ge0:0\le t+kx<1\},\qquad
 c_J(t)=\#\{k\ge0:1\le t+kx<2\}.
\]
Take the natural cost $W_{\rm nat}=2m+N_I+N_J$. Relative to it, the
baseline full-pencil and low-prime rank savings are
\begin{equation}
\begin{aligned}
 w_0&=\max(0,1-A-\gamma,1-B-\gamma,A+B-1-D),\\
 w_{\rm low}&=\max(w_0,1-A-\eta,1-B-\eta).
\end{aligned}                                               \label{eq:7}
\end{equation}
Here $w_0$ is the ordinary rank/potential saving. The term $A+B-1-D$
comes from the entrywise alternative $W_{\rm alt}$ in
Lemma~\ref{lem:residue-copy-rank-insert}. Thus the low-prime cost is
$W_{\rm nat}-mw_{\rm low}(x)$. The conservative branch weight in
\eqref{eq:9} is
\begin{equation}
W_{\rm cons}(p)=
\begin{cases}
4m\ell_p,&p<5\ \text{or}\ p^2\le U+g,\\
\min\{4m,W_{\rm nat}-mw_{\rm low}(p/m)\},
   &5\le p\le g,\ p^2>U+g,\\
\min\{4m,W_A,W_B\},&g<p<m,\\
W_{\rm tail},&m\le p\le U,\\
0,&p>U.
\end{cases}                                                   \label{eq:W-cons-piecewise}
\end{equation}
At a cell boundary $p\mid g$ with $p^2>U+g$, use the second line and
retain that finite prime separately. Away from such boundaries, define
$w_{\rm cons}(x)$ as the limit of $W_{\rm cons}(p)/m$ when $p/m\to x$;
at a boundary take the lower one-sided value. At an actual prime location
we use the displayed finite value $W_{\rm cons}(p)/m$. For the remaining layers $\nu\ge11$, put
$\theta=\gamma/x$. Then
$T_0=A+B=2\theta+\varepsilon_2$ with $0\le\varepsilon_2\le x/4$.
We need only a lower bound for the available saving:
\[
 w_{\rm low}\ge1-\min(A,B)-\min(\gamma,\eta)
       \ge1-\frac {T_0}2-\frac x2.
\]
Consequently the conservative normalized cost satisfies
\begin{equation}
 w_{\rm cons}(x)\le1+3\theta+\frac78x.                    \label{eq:8}
\end{equation}

Set $\delta_0=\rho/11$. The change of variables
$x=\rho/(\nu+1-\theta)$ gives
$\dd x=\rho(\nu+1-\theta)^{-2}\dd\theta$. The midpoint estimate
\[
 \int_0^1\theta\frac{\rho\,\dd\theta}{(\nu+1-\theta)^2}
 \le \frac12\int_0^1\frac{\rho\,\dd\theta}{(\nu+1-\theta)^2}
      +\frac{\rho}{6\nu^3}
\]
and $\sum_{\nu\ge11}\nu^{-3}\le2/441$ give
\begin{equation}
 \int_0^{\delta_0} w_{\rm cons}(x)\,\dd x
 \le\frac52\delta_0+\frac{\rho}{441}+\frac7{16}\delta_0^2.   \label{eq:16a}
\end{equation}
Let $\vartheta$ be Chebyshev's function. For the finite fallback range use a
separate cutoff $\varepsilon>0$:
\begin{equation}
 \sum_{\sqrt{U+g}<p\le\varepsilon m}W_{\rm cons}(p)\log p
 \le4m\,\vartheta(\varepsilon m)
 =(4\varepsilon+o(1))m^2.                                \label{eq:16b}
\end{equation}
The generic branch $p^2\le U+g$ contributes
$O(m^{3/2}\log m)=o(m^2)$. A cell-boundary prime here means a prime
$p\mid g$ with $p^2>U+g$. It is handled by the second line of
\eqref{eq:W-cons-piecewise}; there is at most one such prime, and its
contribution is $O(m\log m)$.

The endpoint weight is
\[
 e_p^{\rm end}=4m\,\ell_p+H_*v_p(a-1)-Nv_p(a).
\]
The local branch is
admissible only when $p\nmid a(a-1)$. Thus we define the coefficientwise
integerizer piecewise by
\begin{equation}
\begin{gathered}
 e_p=\begin{cases}
 e_p^{\rm end},&p\mid a(a-1),\\
 \min\{e_p^{\rm end},W_{\rm cons}(p)\},&p\nmid a(a-1),
 \end{cases}\\
 D_m(a)=\prod_p p^{e_p}.
\end{gathered}                                               \label{eq:9}
\end{equation}
Then $D_m(a)$ is a positive rational number. It satisfies
\begin{equation}
 D_m(a)G_{m,a}(T)\in\mathbb Z[T].                          \label{eq:10}
\end{equation}
Primes dividing $a(a-1)$ take the first branch. So the low-prime unit
hypotheses are never used there. For these primes the $4m\ell_p$ part totals
$O(m\log m)$. The endpoint-divisor terms remain explicitly in
$e_p^{\rm end}$ and in the rate below. We now justify the two exponents. The first
exponent in \eqref{eq:9} follows from the two formal divisors.
For the second, use the separate coefficientwise branch that realizes the
displayed minimum in $W_{\rm cons}(p)$. The ordinary row and column
potentials give the baseline part of $w_0$ and $w_{\rm low}$. The entrywise
alternative is the $W_{\rm alt}$ term in Lemma~\ref{lem:residue-copy-rank-insert}.
The remaining low-prime saving is the coefficientwise rank bound proved
there over the Gauss valuation ring. Thus the low-prime saving is not inferred
from potentials alone. For $g<p<m$, the two valid potential covers are
$W_A,W_B$. For $m\le p\le U$ the displayed tail cover gives
$W_{\rm tail}$. Expanding the determinant after applying any one selected
cover gives the corresponding valuation bound. Division by $b^{mH}$ accounts
for the last term in \eqref{eq:9}. Each branch alone supplies an integerizer,
so taking the pointwise minimum does not combine savings.

\subsection{The limiting arithmetic rates}

The exact low-layer increments for $\nu=1,\ldots,10$, in order, are
\begin{equation}
\begin{aligned}
\rho=5/16:\ &47/6720,\ 69/41888,\ 387/637840,\ 6219367/21864360960,\\
&75/484352,\ 180207/1941459520,\ 317501/5311088640,\\
&18175/441971712,\ 11263/382261760,\ 155/7117824;\\[2pt]
\rho=1/3:\ &25/3024,\ 1/540,\ 4/5859,\ 11/34344,\ 7/40095,\\
&17/161460,\ 2/29295,\ 23/492048,\ 13/389367,\ 29/1175580.
\end{aligned}                                                    \label{eq:11}
\end{equation}
Their sums are
\begin{equation}
 L_{5/16}=\frac{1111075051739028956195167}
 {111864885256841731511777280},\quad
 L_{1/3}=\frac{85930118826854129}{7423773458762901240}.    \label{eq:12}
\end{equation}
Write $I_{\rm high}=\int_\rho^1w_{\rm cons}(x)\,\dd x$ and
$I_{\rm tail}=\int_1^2W_{\rm tail}(x)\,\dd x=2-3\rho^2/4$. Let
$I_{\nu,0}$ and $I_{\nu,1}$ denote the baseline and rank-improved layer
integrals. The baseline profile is a minimum of costs. Its normalized costs
are $4$, $W_{\rm nat}/m-w_0$, $W_A/m$, $W_B/m$, and $W_{\rm tail}/m$.
For $1\le x\le2$ the last one has continuum form
$w_{\rm tail}(x)=4(2-x)-3(2+\rho-2x)_+$.
The ten disjoint increments in \eqref{eq:11} are the exact integrals of
the improvement $w_{\rm low}-w_0$. Each is taken after the pointwise minimum
with the other valid branches. The replay uses the change of variables above,
the midpoint remainder $\rho/(6\nu^3)$, and
$\sum_{\nu\ge11}\nu^{-3}\le2/441$ to include the tail.
Define the replay's explicit upper expression by
\begin{equation}
 I_{\rm upper}=I_{\rm high}+I_{\rm tail}
   +\sum_{\nu=1}^{10}I_{\nu,1}
   +\frac52\delta_0+\frac{\rho}{441}+\frac7{16}\delta_0^2,
 \qquad \delta_0=\rho/11.                                 \label{eq:I-upper}
\end{equation}
The exact profile replay gives
\[
\begin{array}{c|c|c|c}
\rho&I_{\rm high}&I_{\rm tail}&I_{\rm upper}\\ \hline
5/16&98909/53760&1973/1024&4.6343048426478513377\ldots\\
1/3&125/72&23/12&4.5805812506617284878\ldots
\end{array}
\]
The profile replay recomputes the finite rational layer integrals. It verifies
that this upper integral is below the following rational bounds:
\begin{equation}
 \begin{aligned}
 S_-&=\frac{8108653419015955202861020329389}
 {1747888832138152054871520000000}<4.639112779899639,\\
 S_+&=\frac{106301329053850833773025543473}
 {23199292058634066375000000000}<4.582093659805828.
 \end{aligned}                                                \label{eq:13}
\end{equation}
Thus $S_\rho$ is directly a rational upper bound for
$\int_0^2w_{\rm cons}(x)\,\dd x$.

For fixed $\varepsilon>0$, \eqref{eq:16b} controls
$p\le\varepsilon m$. On each fixed cell
$\varepsilon\le p/m\le2$, residue counts are affine up to $O(1)$. We then
apply the prime number theorem with partial summation, in the standard form
of \cite{Davenport}. This gives
\begin{equation}
 \frac1{m^2}\sum_{\alpha m<p<\beta m}
 W_{\rm cons}(p)\log p
 \longrightarrow\int_\alpha^\beta w_{\rm cons}(x)\,\dd x. \label{eq:14}
\end{equation}
The at-most-one cell-boundary prime described above contributes
$O(m\log m)$. Let first $m$ tend to infinity
and then $\varepsilon$ tend to zero. This yields
\begin{equation}
 \limsup_{m\to\infty}\frac{\log D_m(a)}{m^2}
 \le S_\rho+\frac{9\rho^2}{4}\log|a-1|
                  -\frac18\log|a|.                        \label{eq:15}
\end{equation}
Thus, with $q_a=|a|$,
\begin{equation}
\begin{aligned}
 \tau_+(q_a)&=S_++\frac14\log(q_a-1)-\frac18\log q_a,\\
 \tau_-(q_a)&=S_-+\frac{225}{1024}\log(q_a+1)-\frac18\log q_a.
\end{aligned}                                               \label{eq:17}
\end{equation}
The endpoint divisor contributes $H_*/m^2\to9\rho^2/4$ and
$N/m^2\to1/8$. This proves the required arithmetic limsup.

\section{Nonvanishing along prime-power towers}\label{sec:nonvanishing}

The ancillary seed files contain two coefficient arrays. Let $H_{211}$ and
$H_{378}$ be the integer polynomials they determine.
Appendix~\ref{app:seed-data} prints their degrees and complete residue
vectors. The arrays themselves are the exact coefficient data used by the
seed replay. The seed arrays and the interval boxes used later are part of
the proof record. They accompany the paper as ancillary data. Take the seed
pair $(n,h)=(12,4)$ or $(16,5)$, and write $b=a-1$ and $H=3h-3$ at the seed.
We use row normalization, as in the ancillary script. For each seed row
$i\in\{0,\ldots,n-1\}$, let
\[
 s_i=\operatorname{lcd}\{\text{coefficients of }b^H E_{ij}(a):
                         n\le j<2n\},
 \qquad
 N_n(a)=\bigl[s_i b^H E_{ij}(a)\bigr]_{\substack{0\le i<n\\ n\le j<2n}},
\]
and put $S_n=\prod_{i=0}^{n-1}s_i$. Thus $N_n(a)$ is an integer polynomial
matrix and $\det N_n(a)=S_n b^{nH}G_{n,a}(0)$.
Define
\[
\begin{aligned}
K_{12}&=2^{148}3^{74}5^{20}7^{32}11^{37}13^{38}17^{28}
        19^{20}23^4,\\
K_{16}&=2^{154}3^{96}5^{91}7^{32}11^{46}13^{52}17^{51}
        19^{52}23^{36}29^{12}31^4,
\end{aligned}
\]
and let $K_n$ denote the corresponding constant. Put
$f_{12}=25$, $f_{16}=42$, $e_{12}=12$, and $e_{16}=24$. Then the preceding
identity
\begin{equation}
 K_n\det N_n(a)=S_n a^{e_n}(a-1)^{nH-f_n}H(a),
\label{eq:seed-cleared}
\end{equation}
is the exact cleared identity checked by the seed script. Thus the
corresponding determinant quotient has denominator $(a-1)^{f_n}$. The two
sides of \eqref{eq:seed-cleared} have the same polynomial after clearing the
row product. Their
degrees are at most
$\sum_{j=12}^{23}(j+8)=306$ and
$\sum_{j=16}^{31}(j+11)=552$, respectively. These are the column degree
bounds $j+H-1$ summed over the $m=n$ columns. So exact evaluation at
$307$ and $553$ integer points proves the identities. They read
\begin{equation}
 G_{12,a}(0)=\frac{a^{12}H_{211}(a)}
 {K_{12}(a-1)^{25}},\qquad
 G_{16,a}(0)=\frac{a^{24}H_{378}(a)}
 {K_{16}(a-1)^{42}},                                    \label{eq:18}
\end{equation}
On the positive branch write $H=H_{211}$. On the negative branch write
$H=H_{378}$. The residues of $H_{211}$ are nonzero modulo $29$. Those of $H_{378}$ are
nonzero modulo $37$. The residue vectors are printed in Appendix
\ref{app:seed-data}; the ancillary seed replay recomputes them from the
coefficient arrays. Hence
$H(a)\ne0$ for every integer $a$.
For a fixed integer $|a|\ge27$, choose a prime divisor $p$ of
\begin{equation}
 (2n)!\,|a(a-1)H(a)|+1,                                  \label{eq:19}
\end{equation}
where $(n,H)=(12,H_{211})$ on the positive side. On the negative side it
is $(16,H_{378})$. Then $p>2n$ and $p\nmid a(a-1)H(a)K_n$. Every prime
factor of $K_{12}$ is below $24$, and every prime factor of $K_{16}$ is
below $32$. So the seed value in \eqref{eq:18} is a $p$-unit.

Put $Q=p^k$, $(m,g)=(nQ,hQ)$, with $(n,h)=(12,4)$ or $(16,5)$, and work in
$\mathbb Z_{(p)}$. The coefficientwise tower calculation follows.
For $0\le s<Q$ put

\begin{equation}
 B_s(x)=\binom{Qx-s-1}{hQ}.
\end{equation}
Pair the nonmultiples of $Q$ in the numerator window with the corresponding
nonmultiples in $\{1,\ldots,hQ\}$. Each quotient is a $p$-adic unit whose
reduction does not depend on $x$. So the product formula first shows $B_s(x)\in\mathbb Z_{(p)}[x]$.
Partition the $hQ$ numerator factors
according to their residue modulo $Q$. Every residue occurs $h$ times.
More explicitly, pair a numerator factor $Qx-c$ with a denominator factor
$c'$ having the same residue modulo $Q$. Then $c\equiv c'\pmod p$ and
$(Qx-c)/c'\equiv-1\pmod p$. The $h$ factors with $Q\mid c$ give the
remaining product $\binom{x-1}{h}$ after the powers of $Q$ are canceled.
After removing the fixed powers of $p$ contributed by $(hQ)!$, all residue
classes except the one with $t\equiv-s-1\pmod Q$ contribute $p$-adic units
independent of $x$ modulo $p$. The exceptional class contributes exactly
$h$ linear factors in $x$. Thus the reduced polynomial has degree at most
$h$. Lucas' formula gives its values at $x=1,\ldots,h$. At $x=0$ the
negative-index identity gives the same value $(-1)^h$, because
$\binom{-1}{h}=(-1)^h$; on the left $Q=p^k$ is odd since $p>2n$.
The remaining sign is $(-1)^{h(Q-1)}=1$. The $h+1$ values therefore prove the coefficientwise congruence
is the unsigned identity
\begin{equation}
 \binom{Qx-s-1}{hQ}\equiv\binom{x-1}{h}\pmod p.       \label{eq:20}
\end{equation}

Write $i=Qi_0+s_i$ and $j=Qj_0+s_j$, with residues in
$\{0,\ldots,Q-1\}$. Let $P^0_{i_0j_0}$, $E^0=(E^0_{i_0j_0})$ and
$C^0=(C^0_{i_0j_0})$ be the polynomial parts $P_{ij}$ and the matrices
$E(a)$ and $C(a)$ of Section~\ref{sec:determinant} with $(m,g)=(n,h)$.
For a nonzero offset $d$ we use the Taylor series
\[
 U_d(x)=\frac{Q}{Qx+d}=\frac{Q/d}{1+(Q/d)x}.
\]
Since $|d|<2nQ<pQ$, its coefficients are $p$-integral. Put
$d_i=w(-i)/(j-i)$ and $d_j=-w(-j)/(j-i)$. The scaled polynomial part
satisfies the exact identity
\begin{equation}
 Q^2P_{ij}(Qx-s)=B_s(x)^3U_{i-s}(x)U_{j-s}(x)
 -(Qd_i)U_{i-s}(x)-(Qd_j)U_{j-s}(x).                  \label{eq:20a}
\end{equation}
Suppose one offset is zero. The two residue values then agree. So the
numerator in the corresponding $1/x$ term has zero constant coefficient.
Division by $x$ therefore remains integral. Now reduce \eqref{eq:20a}. A nonzero
offset contributes only when it is a multiple of $Q$. Both offsets
contribute precisely when $s_i=s_j=s$. In that case \eqref{eq:20} gives the seed
polynomial part $P^0_{i_0j_0}$. This includes the finite-part division at a
zero offset. After two derivatives we obtain, coefficientwise,
\begin{equation}
 Q^4P_{ij}''(Qx-s)\equiv
 {\bf1}_{s_i=s_j=s}(P^0_{i_0j_0})''(x)\pmod p.             \label{eq:21}
\end{equation}

The scaled values in \eqref{eq:21} have integral Newton coefficients. So Newton
interpolation and
\[
 \sum_{r\ge1}z^r\binom{r-1}{t}=
 \frac{z^{t+1}}{(1-z)^{t+1}}
\]
sum \eqref{eq:21} coefficientwise. Split $r=Qx-s$. A matching residue then
contributes the factor $a^s$ and the denominator $1-a^{-Q}$. This
denominator is a $p$-unit because $p\nmid a(a-1)$, so
$a^Q\equiv a\not\equiv1\pmod p$. The
zero-offset finite part is included before this splitting. Because $h<p$,
the $p$-parts of the factors in $(hQ)!$ cancel coefficientwise against the
corresponding numerator blocks. For the harmonic term,
\[
 Q^3H_{Qu+s}(a)\equiv a^sH_u(a)\pmod p,
\]
because only the multiples $t=Qv$ survive and $a^Q=a$. We combine the
polynomial part, the harmonic term, and the $T$ coefficient. This gives the
all-$k$ congruences
\begin{equation}
 Q^4E_{ij}\equiv {\bf1}_{s_i=s_j}a^sE^0_{i_0j_0},\qquad
 QC_{ij}\equiv {\bf1}_{s_i=s_j}a^sC^0_{i_0j_0}\pmod p.       \label{eq:22}
\end{equation}

Define
\[
 {\mathcal P}_k(X)=Q^{4m}G_{m,a}(X/Q^3).
\]
Group rows and columns by their residues modulo $Q$. The block permutation
signs cancel. The $Q$ diagonal blocks are
$a^s(E^0+XC^0)$. Their scalar is
$a^{nQ(Q-1)/2}=1$ in $\mathbb F_p$, since $n$ is even and
$a^{Q-1}=1$. Thus, for every $k$,
\begin{equation}
 {\mathcal P}_k(X)\in\mathbb Z_{(p)}[X],\qquad
 {\mathcal P}_k(X)\equiv G_{n,a}(X)^Q\pmod p.                  \label{eq:23}
\end{equation}
Let $\xi$ be fixed and rational. Choose $k$ with
$3k+v_p(\xi)\ge1$ (with $v_p(0)=+\infty$). Evaluate at
$X=Q^3\xi$. For all sufficiently large $k$ this gives the unscaled identity
\[
 Q^{4m}G_{m,a}(\xi)={\mathcal P}_k(Q^3\xi)
 \equiv G_{n,a}(Q^3\xi)^Q
 \equiv G_{n,a}(0)^Q\not\equiv0\pmod p.
\]
The seed choice in \eqref{eq:19} makes $G_{n,a}(0)$ a $p$-unit. Since the
left side would be zero if $G_{m,a}(\xi)=0$, this unit congruence gives the
required contradiction. The argument uses the displayed congruence before
any cancellation of the factor $Q^{4m}$.


\section{Positive analytic estimate}\label{sec:positive}

Let $m=12Q$, $g=4Q$, $\rho=1/3$, and $q=|a|>1$, $\lambda=\log q$. Define
\[
 h_\rho(t)=t\log t-(t-\rho)\log(t-\rho)-\rho\log\rho,\quad
 B(z)=\int_0^2\log|z+x|\,\dd x\quad(z\in\mathbb C\setminus[-2,0]).
\]
For real $t\ge\rho$, this complex-plane definition is the real function
$B(t)$ used in the energy below.
On probability densities $\mu=f(t)\dd t$ with $0\le f\le3$, support in
$[\rho,\infty)$, and finite first moment, put
\[
 L(\mu,\nu)=-\iint\log|t-s|\,\dd\mu(t)\dd\nu(s).
\]
\[
 {\mathcal E}^+_\lambda(\mu)=\frac32+\int(B-3h_\rho+\lambda t)\,\dd\mu
 +L(\mu,\mu),
 \quad A_+(\lambda)=\inf_\mu{\mathcal E}^+_\lambda(\mu).
\]

\begin{proposition}\label{prop:positive}
For the positive tower,
\begin{equation}
 \limsup_{Q\to\infty}(12Q)^{-2}
 \log|\det M(a)|\le-A_+(\log|a|).                           \label{eq:24}
\end{equation}
Moreover $A_+$ is finite and concave, and
\begin{equation}
 A_+(\lambda_2)-A_+(\lambda_1)
 \ge\frac12(\lambda_2-\lambda_1)\quad
 (\lambda_2\ge\lambda_1).                                  \label{eq:25}
\end{equation}
\end{proposition}

\begin{proof}
For every center $z$ used below, the circle $|u-z|=1/4$ avoids the
negative-integer poles. Cauchy's formula gives
\begin{equation}
 F_{ij}''(z)=\frac{2!}{2\pi i}
 \int_{|u-z|=1/4}\frac{F_{ij}(u)}{(u-z)^3}\,\dd u.
 \label{eq:cauchy-fpp}
\end{equation}
Thus the two derivatives can be represented before the determinant summation.
The representation also makes the triple zero of $w$ and the two derivative
orders explicit in the occupancy estimate below.
Expand the finite sums $r\le R$ and apply Andreief's identity
\cite{Andreief} to the two Cauchy determinants. Then let $R$ tend to
infinity. Absolute convergence holds because the differentiated rational
summands have polynomial growth and $q^{-\sum r_\ell}$ is geometric. The
Cauchy double alternant gives the standard formula
\cite[Eq.~(2.7)]{Krattenthaler}
\[
 H(z_1,\ldots,z_m)=
 \Delta(I)\Delta(J)\Delta(z)^2
 \prod_{\ell=1}^m\frac{w(z_\ell)}
 {\prod_{k=0}^{2m-1}(z_\ell+k)}.
\]
Consider a nonzero term. The triple zeros of $w$ and the two derivatives
give it occupancy at most three at every integer center. Suppose $k$
variables share a center. Their squared Vandermonde vanishes to order
$k(k-1)$ on that diagonal. The derivatives have total order $2k$. Since
$k(k-1)>2k$ for $k\ge4$, every such term vanishes. Use Cauchy circles of
radius $1/4$ independently in each variable.

Use the midpoint convention $s_r=r+1/2$ and
$\zeta_r=s_r/m$. Let $\mu_m$ have density equal to the center occupancy on
$[r/m,(r+1)/m)$. In cells a fixed positive distance from the support endpoint,
replacing $s_r/m$ by either endpoint changes each displayed Riemann sum by
$O(1/m)$ per particle. The bounded collection of cells meeting the endpoint
has $h'_{1/3}(t)=\log(t/(t-1/3))$ and contributes
$O(\log m/m)$ per particle. The monotone tail estimates handle the remaining
cells. Summed over all particles, these errors are $O(m\log m)$.
The fixed Vandermonde, moving Vandermonde, filter, and denominator
estimates hold uniformly in the admissible center tuples. With this
convention they are
\begin{equation}
\begin{aligned}
 \log(\Delta(I)\Delta(J))
   &=m^2\log m-\frac32m^2+O(m\log m),\\
 \log|\Delta(z)^2|
   &\le m(m-1)\log m-m^2L(\mu_m,\mu_m)+O(m\log m),\\
 \sum_\ell\log|w(z_\ell)|
   &\le3m^2\int h_\rho\,\dd\mu_m+O(m\log m),\\
 \sum_{\ell,k}\log|z_\ell+k|
   &\ge2m^2\log m+m^2\int B\,\dd\mu_m-O(m\log m).
\end{aligned}                                             \label{eq:26}
\end{equation}
The filter estimate is the entropy inequality for $\binom{r-1}{g}$. It is
followed by the increasing Riemann sum for $h_\rho$. For the moving
Vandermonde, let $c_r\le3$ be the occupancy of cell $r$. Remove, for each
particle, pairs with distinct particles in cells at distance $|r-s|<3$.
Occupancy at most three leaves at most fourteen such non-diagonal pairs per
particle. There are two in its own cell and three in each of the four
adjacent cells.
The same-cell integral is $-\log m-3/2$, and the removed pairs contribute
$O(m\log m)$. For the remaining pairs, if $d=|r-s|\ge3$, two particles in these cells are at distance at most
$(d+1)/m$. Integrating over the two cells gives the comparison
\[
 \log\frac{d+1}{m}\le
 m^2\!\int_{r/m}^{(r+1)/m}\!\int_{s/m}^{(s+1)/m}
       \log|x-y|\,\dd x\dd y+O(d^{-1}).
\] Summing this inequality with
weights $c_rc_s$ yields the moving-Vandermonde line in \eqref{eq:26}.
The constants in the $O$ terms can be chosen independently of the tuple of
centers. For $c_r\le3$ and $\sum_rc_r=m$,
\[
 \begin{aligned}
 &\left(\sum_r c_r(c_r-1)+
       \sum_r\sum_{0<|s-r|<3}c_rc_s\right)(1+\log(m+1))\\[-2pt]
 &\hspace{4em}\le 14m(1+\log(m+1)),\\
 &\sum_r\sum_{|s-r|\ge3}\frac{c_rc_s}{|r-s|}
       \ll m\log(m+1).
 \end{aligned}
\]
The first bound handles the deleted same-cell and neighboring-cell pairs.
The second sums the $O(d^{-1})$ comparison error, using at most six particles
in each fixed positive distance shell. The midpoint, endpoint, and
Cauchy-circle replacements satisfy the same uniform estimate
\[
 |R_{\rm edge}|+|R_{\rm circle}|\le C m\log(m+1)
\]
for an absolute $C$ and every admissible center tuple. In particular the
singular first-cell derivative contributes $O(\log m/m)$ per particle before
the outer factor $m$. So this bound includes it.
The denominator line is the increasing Riemann sum for $B$. For $\rho=1/3$
the changed uniform inequalities are
\[
\begin{aligned}
 \log|w(m\zeta)|&\le3m\int_0^{1/3}\log|\zeta-x|\,\dd x
 +3m(\rho-\rho\log\rho)+O(\log m),\\
 \sum_{k=0}^{2m-1}\log|\zeta+k/m|&\ge mB(\zeta)-\log7.
\end{aligned}
\]
The support is $[1/3,\infty)$. For one particle the denominator-circle
loss is
\[
 \sum_{k=0}^{2m-1}-\log\left(1-\frac1{4(r+k)}\right)<2,
 \qquad r+k\ge m/3.
\]
Thus the total circle loss is $<2m$. The denominator line loses $\log7$ per
particle, and hence $m\log7$ in total. Both losses are $O(1)$ per particle,
and $O(m)$ in total.
The fixed and moving Vandermonde bounds use the same near-pair correction
above. The remaining terms are the Riemann sums of the two displayed
inequalities over the cells. Their errors are $O(\log m)$ per particle, so
the total error is $O(m\log m)$. None of these terms alters the quadratic
rate.

Reserve $\exp(-\sum r_\ell/m)$ from the geometric weight and put
$\beta=\lambda-1/m$. Equation~\eqref{eq:26}, Cauchy's estimate, and
cancellation of all $m^2\log m$ terms give each nonzero tuple the bound
\[
 \exp\{-m^2{\mathcal E}^+_\beta(\mu_m)+O(m\log m)
                 -\tfrac1m\sum_\ell r_\ell\}.
\]
The unrestricted tail satisfies
\begin{equation}
 \frac1{m!}\sum_{r_1,\ldots,r_m\ge1}
 e^{-\sum r_\ell/m}
 \le\frac1{m!}(e^{1/m}-1)^{-m}\le e^m.                    \label{eq:27}
\end{equation}
Take the supremum over center tuples, sum \eqref{eq:27}, and then let
$Q\to\infty$. This proves the bound with $A_+(\lambda-1/m)$. The energy is
finite on a compact reference density. It is bounded below by the
confining linear term and is affine in $\lambda$. Hence it is concave and
continuous. Letting $m\to\infty$ gives \eqref{eq:24}.

Finally, we use the bathtub inequality for the density cap $3$. It shows that
every admissible probability density supported on $[\rho,\infty)$ has
first moment at least $\rho+1/6=1/2$. Taking the difference of the two
affine functionals in $\lambda$ proves \eqref{eq:25}. The finite cutoff is
removed exactly as in \cite[Lemma~6.1]{ZhB}. Here the summand is bounded
by the geometric factor in \eqref{eq:27}.
\qedhere
\end{proof}

\section{Negative signed-contour estimate}\label{sec:negative}

Let $m=16Q$, $g=5Q$, $\rho=5/16$, $q=|a|>1$, and $\lambda=\log q$. Put
\[
 H_\rho(z)=\int_0^\rho\log|z-x|\,\dd x-\rho\log\rho+\rho,\quad
 B(z)=\int_0^2\log|z+x|\,\dd x,
\]
\[
 V_\lambda(z)=B(z)-3H_\rho(z)+\lambda\operatorname{Re} z+\pi|\operatorname{Im} z|.
\]
In this section set
\[
 L(\mu,\nu)=-\iint\log|z-w|\,\dd\mu(z)\dd\nu(w).
\]
Let
\[
 \Gamma_m(t)=g+\frac12+4|t|+it,\qquad t\in\mathbb R,
\]
oriented by increasing $t$. Let $C$ be its limiting scaled contour
$z=\rho+4|y|+iy$. For probability measures $\mu$ on $C$ define
\[
 {\mathcal E}^-_\lambda(\mu)=\frac32+\int_CV_\lambda\,\dd\mu
 +L(\mu,\mu),
 \qquad A_-(\lambda)=\inf_\mu{\mathcal E}^-_\lambda(\mu).
\]

\begin{proposition}\label{prop:negative}
For the negative tower,
\begin{equation}
 \limsup_{Q\to\infty}(16Q)^{-2}\log|\det M(-q)|
 \le-A_-(\log q).                                         \label{eq:28}
\end{equation}
The infimum is over all probability measures of finite first moment and
finite logarithmic energy; there is no density-cap assumption. Further,
\begin{equation}
 A_-(\lambda_2)-A_-(\lambda_1)
 \ge\frac5{16}(\lambda_2-\lambda_1)\quad
 (\lambda_2\ge\lambda_1).                                 \label{eq:29}
\end{equation}
\end{proposition}

\begin{proof}
Write $c=g+1/2$. Consider
$F_{ij}(s)=w(s)/((s+i)(s+j))$. Close the contour formed by the two rays
and the downward vertical segment at $\operatorname{Re} s=R$. Let $R$ tend
to infinity through half-integers. The boundary is clockwise. The residue of
$q^{-s}F_{ij}''(s)/\sin(\pi s)$ at a positive integer $r$ is
$(-1)^rq^{-r}F_{ij}''(r)/\pi$. The residues with $r\le g$ vanish because
$w$ has a triple zero. The vertical segment is $O(q^{-R}R^C)$ and tends to
zero. The two rays converge exponentially. Therefore, with orientation,
the exact signed identity is
\begin{equation}
 M_{ij}(-q)=-\frac1{2i}\int_{\Gamma_m}
       \frac{q^{-s}F_{ij}''(s)}{\sin(\pi s)}\,\dd s.         \label{eq:30}
\end{equation}
Set $W_{I,J}=\Delta(I)\Delta(J)$ and
\[
 P(s_1,\ldots,s_m)=
 \frac{\prod_\ell w(s_\ell)\Delta(s)^2}
 {\prod_{\ell=1}^{m}\prod_{k=0}^{2m-1}(s_\ell+k)}.
\]
Finite determinant expansion followed by absolute convergence gives the
differentiated Andreief identity
\begin{equation}
 \det M(-q)=\frac{W_{I,J}}{m!}\left(-\frac1{2i}\right)^m
 \int_{\Gamma_m^m}\prod_\ell\frac{q^{-s_\ell}}{\sin(\pi s_\ell)}
 \left(\partial_1^2\cdots\partial_m^2P\right)
 \prod_\ell\dd s_\ell.                                   \label{eq:31}
\end{equation}
The derivatives act on $P$ only.

At $s=\Gamma_m(t)$,
\begin{equation}
 |\sin(\pi(c+4|t|+it))|^{-1}\le C_{\sin}e^{-\pi|t|}.             \label{eq:32}
\end{equation}
Here $C_{\sin}$ is an absolute constant. The squared modulus is
$\cos^2(4\pi|t|)+\sinh^2(\pi t)$. On $|t|\le1/16$ the cosine is at least
$1/\sqrt2$ in absolute value. On $1/16\le|t|\le1$ the hyperbolic sine has a
fixed positive lower bound. For $|t|\ge1$ it gives the exponential
bound. Cauchy circles of radius $1/8$ give
\begin{equation}
 |\partial_1^2\cdots\partial_m^2P(s)|
 \le128^m\sup_{|u_\ell-s_\ell|\le1/8}|P(u)|.               \label{eq:33}
\end{equation}
We shift only the majorant $q^{-\operatorname{Re}s}e^{-\pi|\operatorname{Im}s|}$
of the base and sine factors in \eqref{eq:32}. On each such circle,
replacing the center $s_\ell$ by a supremum point $u_\ell$ changes this
majorant by at most $\exp((\lambda+\pi)/8)$. The $m$ shifts therefore cost
$(\lambda+\pi)m/8=O_\lambda(m)$ and are absorbed in the linear error term.
Use the scaled coordinate $\zeta=s/m$ throughout. Thus the contour centre
has $\zeta=z(y)+1/(2m)$, and the Cauchy circles give the tube
\[
 {\mathcal S}_m=\{z(y)+1/(2m)+e:\ |e|\le1/(8m)\},\quad
 z(y)=\rho+4|y|+iy.
\]
The endpoint and midpoint Riemann sums have total error
$O(\log m+\log(1+|\zeta|))$ per particle. In particular
$\operatorname{dist}({\mathcal S}_m,C)\le5/(8m)$. The filter and
denominator bounds on this tube are
\begin{equation}
\begin{aligned}
 \log|w(m\zeta)|\le3mH_\rho(\zeta)
 +O\!\left(\log m+\log(1+|\zeta|)\right),\\
 \sum_{k=0}^{2m-1}\log|\zeta+k/m|
 \ge mB(\zeta)-O\!\left(\log m+\log(1+|\zeta|)\right).    \label{eq:34}
\end{aligned}
\end{equation}
The filter bound follows from the decreasing Riemann sum. The denominator
bound follows from the increasing one.
The normalization in $H_\rho$ uses $\log(g!)\ge g\log g-g$. The fixed
Vandermonde satisfies
\[
 \log W_{I,J}=m(m-1)\log m-\frac32m^2+O(m\log m),
\]
and the powers of $m$ cancel against the moving Vandermonde and the $2m$
denominator factors. The resulting integrand is bounded by
\begin{multline}
 \exp\left\{-\frac32m^2-m\sum_\ell V_\lambda(\zeta_\ell)
 +\sum_{\ell\ne j}\log|\zeta_\ell-\zeta_j|\right.\\
 \left.{}+C_2\sum_\ell\log(1+|\zeta_\ell|)+O(m\log m)\right\}.  \label{eq:35}
\end{multline}
Here $C_2=C_2(\lambda)$ is independent of $m$ and of the configuration. It
collects the per-particle errors of \eqref{eq:34}; its $\lambda$ dependence
also absorbs the preceding circle-shift cost.

At a contour center $s=m\zeta$ the base and sine factors satisfy
\[
 q^{-\operatorname{Re} s}e^{-\pi|\operatorname{Im} s|}
 \le \exp\{-m(\lambda\operatorname{Re}\zeta+\pi|\operatorname{Im}\zeta|)\}.
\]
For $L\ge0$
 put $\ell_L(d)=\max(\log d,-L)$, with $\ell_L(0)=-L$. If
$\mu_m=m^{-1}\sum_\ell\delta_{\zeta_\ell}$, then for every $L$
\begin{equation}
 \sum_{\ell\ne j}\log|\zeta_\ell-\zeta_j|
 \le m^2\iint\ell_L(|z-w|)\,\dd\mu_m(z)\dd\mu_m(w)+mL. \label{eq:35a}
\end{equation}
Subtract and add $\varepsilon|\operatorname{Im}\zeta_\ell|$ in \eqref{eq:35}. The bound for the full
configuration then has three parts. They are the corresponding truncated
energy plus $mL+O(m\log(m+1))$, the term $C_2\sum_\ell\log(1+|\zeta_\ell|)$, and the
factor $\exp(-m\varepsilon\sum_\ell|y_\ell|)$. On the tube
$|\zeta|\le1+5|y|$, so $\log(1+|\zeta|)\le\log2+\frac52|y|$. For
$m\ge5C_2/\varepsilon$ we get
$C_2\log(1+|\zeta_\ell|)\le C_2\log2+\frac12m\varepsilon|y_\ell|$. Half of
the factor absorbs these terms, at the cost $e^{O(m)}$. Since
$|\dd s|=m\sqrt{17}\,\dd y$ on the two rays, the remaining factor
$\exp(-\frac12m\varepsilon\sum_\ell|y_\ell|)$ has integral at most
$(4\sqrt{17}/\varepsilon)^m$. Taking first $m\to\infty$, then $L\to\infty$,
and finally $\varepsilon\downarrow0$ proves \eqref{eq:28} using
\eqref{eq:37}--\eqref{eq:38a}.


For $0<\varepsilon\le\pi/2$ put
\[
 V_{\lambda,\varepsilon}(z)=V_\lambda(z)-\varepsilon|\operatorname{Im} z|,\qquad
 K_{\varepsilon,L}(z,w)=\frac{V_{\lambda,\varepsilon}(z)+V_{\lambda,\varepsilon}(w)}2
                         -\ell_L(|z-w|),
\]
and let ${\mathcal E}_{\varepsilon,L}$ be $3/2$ plus the $K_{\varepsilon,L}$
pair integral. Work in the closed region
\[
 D=\{z:\rho-1/32\le\operatorname{Re} z\le\rho+4|\operatorname{Im} z|+1\},
\]
which contains $C$ and every $\mathcal S_m$ for $m\ge16$. There are
constants $c>0,C_0$ such that
\begin{equation}
 K_{\varepsilon,L}(z,w)\ge c(|z|+|w|)-C_0\qquad(z,w\in D).       \label{eq:36}
\end{equation}
These constants are independent of $L$ and $\varepsilon$. The same
elementary logarithmic bounds used for $C$ apply on the fixed strip
$\rho-1/32\le\operatorname{Re}z\le\rho$. Their constants are enlarged
there.
They give
$V_{\lambda,\varepsilon}(z)\ge(\pi/2)|\operatorname{Im} z|-3\rho\log(1+|z|)-C_1$.
Here $C_1$ is independent of $L$ and $\varepsilon$.
On $D$, $|z|\le\rho+1+5|\operatorname{Im} z|$, and
$\ell_L(|z-w|)\le\log(1+|z|)+\log(1+|w|)$. This proves \eqref{eq:36}.

Let $a_{\varepsilon,L,m}$ be the infimum of ${\mathcal E}_{\varepsilon,L}$ over
probability measures on $\mathcal S_m$. Let $a_{\varepsilon,L}$ be the
corresponding infimum on $C$. Choose a fixed compactly supported reference
measure on a finite piece of $C$. Push it by the real shift $1/(2m)$.
Its support is contained in $\mathcal S_m$, so it gives a uniform upper
bound for near-minimizers of $a_{\varepsilon,L,m}$. Together
with the coercivity in \eqref{eq:36}, this gives tightness and a uniform
first-moment bound. From any weakly convergent subsequence, the distance
bound $\operatorname{dist}({\mathcal S}_m,C)\le5/(8m)$ puts the limit on
$C$. For fixed $L$, $K_{\varepsilon,L}$ is bounded below and continuous. The
product measures converge weakly. We truncate the kernel above and then use
monotone convergence. This proves the lower-semicontinuity inequality
\begin{equation}
 \liminf_{m\to\infty}a_{\varepsilon,L,m}\ge a_{\varepsilon,L}.   \label{eq:37}
\end{equation}
As $L$ increases, $\ell_L$ decreases and $K_{\varepsilon,L}$ increases to the
full kernel. Hence $a_{\varepsilon,L}$ is nondecreasing and has a limit no
larger than the full infimum $A_{-,\varepsilon}(\lambda)$. For the reverse
inequality, take near-minimizers for $a_{\varepsilon,L}$ and use
\eqref{eq:36} for tightness. Pass to a weak limit under each fixed
truncation $K_{\varepsilon,L_0}$, then let $L_0$ increase. This identifies
the limit of the infima with the full infimum.

It remains to remove the auxiliary $\varepsilon$. Let
$\varepsilon_j\downarrow0$ and choose $\mu_j$ on $C$ with
${\mathcal E}_{\varepsilon_j}(\mu_j)\le A_{-,\varepsilon_j}(\lambda)+1/j$.
The same reference measure and \eqref{eq:36} give tightness and a uniform
first-moment bound. After passing to a subsequence, $\mu_j\Rightarrow\mu$.
Fix $L$ and a clipping level $T$, and set
$K_{j,L}^T=\min(K_{\varepsilon_j,L},T)$. Choose a continuous cutoff
$\chi_R$ on $D$. For $R>0$ define the compact exhaustion
$D_R=D\cap\{|z|\le R\}$. Choose $\chi_R$ equal to $1$ on $D_R$ and
supported in $D_{R+1}$, with $0\le\chi_R\le1$. On the compact support of
$\chi_R(z)\chi_R(w)$, the kernels $K_{j,L}^T$ converge uniformly to
$K_{0,L}^T$. The coercive lower bound \eqref{eq:36} controls the complement,
and tightness makes its contribution at least $-2\eta(T+C_0)$ for $R$ large.
Weak convergence applied to the continuous cutoff therefore gives
\[
 \liminf_{j\to\infty}\iint K_{\varepsilon_j,L}\,\dd\mu_j\dd\mu_j
 \ge \iint K_{0,L}^T\,\dd\mu\dd\mu-2\eta(T+C_0).
\]
First let $j\to\infty$, then $R\to\infty$ and $\eta\downarrow0$, and
finally $T\to\infty$. This yields the varying-kernel lower bound at fixed
$L$. Letting $L\to\infty$ gives
\begin{equation}
 \lim_{\varepsilon\downarrow0}A_{-,\varepsilon}(\lambda)=A_-(\lambda). \label{eq:38a}
\end{equation}
We turn to the zero-mass energy inequality used in endpoint certificates.
Let $\mu$ and $\sigma$ be probability measures of equal mass, with $\sigma$
compactly supported. Assume that $\mu$ has finite first moment and finite
logarithmic self-energy. Assume also that the cross-pairing with $\sigma$ is
finite.
Put $\nu=\mu-\sigma$. Restrict both
measures to a disk of radius $R$ and normalize. This gives probability
measures $\mu_R$ and ${\sigma}_R$. Put $\nu_R=\mu_R-{\sigma}_R$. For all
large $R$, ${\sigma}_R=\sigma$, so $\nu_R$ is compactly supported and has
zero mass.
The compact zero-mass energy inequality
\begin{equation}
 L(\nu_R,\nu_R)\ge0                                      \label{eq:38}
\end{equation}
is the direct compact case of \cite[Lemma~I.1.8]{SaffTotik}. We use only
this inequality; no Fourier limiting identity is needed.
The positive part of $\log|z-w|$ is bounded by
$\log(1+|z|)+\log(1+|w|)$. The negative part is controlled by the finite
self-energy. Also $U^\sigma(z)=O(\log(1+|z|))$. Consequently
$L(\mu_R,\mu_R)\to L(\mu,\mu)$ and
$L(\mu_R,\sigma)\to L(\mu,\sigma)$, so
$L(\nu_R,\nu_R)\to L(\nu,\nu)$. We pass only this pairing limit in
\eqref{eq:38}:
\begin{equation}
 L(\mu-\sigma,\mu-\sigma)\ge0.                         \label{eq:38b}
\end{equation}
The bounded compact densities of both endpoint certificates satisfy these
hypotheses.

Since $\operatorname{Re} z\ge\rho=5/16$ on $C$, subtracting the two affine
functionals in $\lambda$ gives \eqref{eq:29}.
\qedhere
\end{proof}

\section{Endpoint certificates and ray comparison}\label{sec:endpoints}

We first state the certificate inequalities and then record their
numerical values. Consider the positive problem first. Let
$\sigma=f_\sigma(t)\dd t$ be the piecewise-constant density given by the
$256$-cell table in
Appendix~\ref{app:positive-data}. Put
\[
\begin{aligned}
 U^\sigma(t)&=\int\log|t-s|\,\dd\sigma(s),\\
 V_+(t)&=B(t)-3h_{1/3}(t)+\lambda t,\\
 \Psi_+(t)&=V_+(t)-2U^\sigma(t).
\end{aligned}
\]
For every real $\Lambda$,
\begin{equation}
 A_+(\lambda)\ge \frac32-L(\sigma,\sigma)+\Lambda
       -3\int_{1/3}^{\infty}(\Lambda-\Psi_+(t))_+\,\dd t.
                                                               \label{eq:41}
\end{equation}
Take an admissible density $f\le3$. The pointwise inequality
$(\Psi_+-\Lambda)f\ge-3(\Lambda-\Psi_+)_+$ and
$L(f\dd t-\sigma,f\dd t-\sigma)\ge0$ give \eqref{eq:41}.
Explicitly,
\[
L(\mu,\mu)=L(\mu-\sigma,\mu-\sigma)-2\int U^\sigma\,\dd\mu-L(\sigma,\sigma),
\]
which fixes the sign in \eqref{eq:41}. The zero-mass inequality is the
one-dimensional logarithmic-energy inequality; it is the real-line special
case of \eqref{eq:38b}. On each printed cell, the convex
terms of $\Psi_+$ are bounded by their midpoint tangent. The concave terms
are bounded by their endpoint chord. The tail is bounded by
\[
 \Psi_+(t)\ge B(t)-3h_{1/3}(t)+\lambda t-2\log(t-1/3),
\]
whose derivative is at least $\lambda-(3\rho+2)/(t-\rho)$. The interval
check verifies that this bound is increasing beyond the last cell.

Consider the negative problem. Let $\sigma$ be the positive rational
probability measure on the two rays in Appendix~\ref{app:negative-data}. The compact
certificate uses the
interval boxes in the ancillary files. On the whole contour it verifies
\begin{equation}
\Psi_-(z):=V_\lambda(z)-2U^\sigma(z)\ge\Lambda_-
 \quad(z\in C),                                          \label{eq:42}
\end{equation}
The contour, $V_\lambda$, and the certificate measure are invariant under
complex conjugation. Thus the interval replay on the displayed $+$ ray gives
the same bound on the $-$ ray. The replay also evaluates
$L(\sigma,\sigma)$ with the exact pair kernel, including the displayed
minus sign. This
evaluation includes the finite part at the cusp. Since every competitor
has mass one, \eqref{eq:42} and
$L(\mu-\sigma,\mu-\sigma)\ge0$ give
\begin{equation}
 A_-(\lambda)\ge\frac32-L(\sigma,\sigma)+\Lambda_- .     \label{eq:43}
\end{equation}
For the two endpoint values, the replay computes $\Lambda_-$ as the minimum
of $V_\lambda-2U^\sigma$ over the directed interval boxes and rounds that
minimum downward:
\begin{equation}
 \Lambda_-^{(27)}=5.0771336769,\qquad
 \Lambda_-^{(28)}=5.1031720708.                            \label{eq:lambda-minus}
\end{equation}
The interval check subdivides compact ray pieces. It bounds the infinite
tail by the monotone derivative of $V_\lambda-2U^\sigma$. It rounds the
resulting lower bounds downward. Thus these are directed certificates.
They do not rest on any numerical optimization.

The positive endpoint levels used in the two positive checks are
\[
 \Lambda_{27}=\frac{1364590351}{268435456},\qquad
 \Lambda_{28}=\frac{687124203}{134217728}.
\]
Each printed positive cell is subdivided into $32$ equal subcells. Write the
subcells as $[t_i^-,t_i^+]$ and put $h_i=t_i^+-t_i^-$. Let $f_i$ be the
constant density on that subcell, and let $\mu_i=f_i h_i$ be its mass. With
$\Phi(x)=x^2(\log|x|/2-3/4)$, using the real logarithm and the continuous
value $\Phi(0)=0$, the exact pair term is
\[
 L(\sigma,\sigma)=-\sum_{i\ge j}\frac{\mu_i\mu_j}{h_i h_j}
 (2-\mathbf1_{i=j})
 [\Phi(t_i^+-t_j^-)-\Phi(t_i^--t_j^-)-\Phi(t_i^+-t_j^+)+\Phi(t_i^--t_j^+)].
\]
After the last cell the check uses the increasing tail
$V_+(t)-2\log(t-1/3)$.

For the negative table, the branch directions are
$A=(4+i)/\sqrt{17}$ and $\overline A=(4-i)/\sqrt{17}$. A cell count satisfies
$c_j\ge0$ and $\sum_jc_j=d$ for each printed column. It has mass
$c_j/(2d)$ on each branch and density
$c_j/(2d\,\Delta u_j)$. Here $d$ is the printed denominator and
$\Delta u_j$ is the cell width. For the cross-ray primitive set
$\beta=-\overline A=(-4+i)/\sqrt{17}$, so
$Au+\beta v=z_+(u)-z_-(v)$. Use the principal analytic branch
$\operatorname{Log} z$
on $\mathbb C\setminus(-\infty,0]$ and define
$\Xi(z)=z^2(\operatorname{Log} z/2-3/4)$ there, with $\Xi(0)=0$ by continuity. Take a cell
$[t^-,t^+]$ on the first
branch and a cell $[s^-,s^+]$ on the second. The opposite branch uses
the signed contribution
\[
-\operatorname{Re}\!\left(\frac{\Xi(At^++\beta s^+)-\Xi(At^-+\beta s^+)
 -\Xi(At^++\beta s^-)+\Xi(At^-+\beta s^-)}{A\beta}\right),
\]
with these density factors. The interval lower bounds are
$4.97200113\ldots$ and $4.99000040\ldots$ before downward rounding. The
largest support endpoint is $u_{\max}=238521/262144$. To justify the tail
bound, write $A=(4+i)/\sqrt{17}$ and $z_+(u)=\rho+Au$. For $t_0\ge0$,
\[
 \operatorname{Re}\frac{A}{t_0+Au}
 =\frac{u+4t_0/\sqrt{17}}
        {u^2+8t_0 u/\sqrt{17}+t_0^2},
 \qquad 0\le\operatorname{Re}\frac{A}{t_0+Au}\le\frac1u.
\]
The same formula holds on the conjugate ray. The logarithmic terms of the function $B(z)$ have
$t_0=\rho+x\ge0$. Those of $H_\rho(z)$ have $t_0=\rho-x\ge0$. On the tail $u\ge3$, beyond the last cell edge, the
formula gives
\[
 \frac{\dd}{\dd u}\bigl(B(z_+(u))-3H_\rho(z_+(u))\bigr)
 \ge-\frac{3\rho}{u}\ge-\rho.
\]
Take a support point $z_\pm(v)$ with $0\le v\le u_{\max}$. Translate the
origin along the ray. The same identity then gives
\[
 \left|\frac{\dd}{\dd u}\log|z_+(u)-z_\pm(v)|\right|
 \le\frac1{u-u_{\max}}\qquad(u>u_{\max}).
\]
The certificate measure is positive with total mass one, so the contribution of
$-2U^\sigma$ is at least $-2/(u-u_{\max})$. Therefore, on either ray,
the tail derivative is bounded by
\[
 \frac{4\lambda+\pi}{\sqrt{17}}-\rho-\frac{2}{u-u_{\max}}>0.
\]

The exact rational and directed interval checks use the tables in
Appendices~\ref{app:positive-data} and~\ref{app:negative-data}.
They give the downward-rounded energy values
\begin{equation}
\begin{array}{c|cc}
 &q=27&q=28\\ \hline
 A_+(\log q)&\ge4.9972574688&\ge5.0224702721\\
 A_-(\log q)&\ge4.9720000000&\ge4.9900000000
\end{array}                                                \label{eq:44}
\end{equation}
and the arithmetic costs \eqref{eq:17}. The costs use rational
enclosures of the logarithms. Each margin below
is energy minus arithmetic cost. The resulting strict margins are
\begin{equation}
\begin{array}{c|cc}
 &q=27&q=28\\ \hline
 \text{positive}&0.0126192827393433594&0.0329429595649909584\\
 \text{negative}&0.0126929857831087892&0.0275284462287496169 .
\end{array}                                                \label{eq:margins}
\end{equation}
For $q\ge27$,
\[
 \frac{\dd\tau_+}{\dd\log q}
 =\frac{q}{4(q-1)}-\frac18\le\frac7{52}<\frac{27}{200}.
\]
Together with \eqref{eq:25}, the positive margin increases with slope at
least $73/200$. On the negative side,
\[
 \frac{\dd\tau_-}{\dd\log q}
 =\frac{225}{1024}\frac q{q+1}-\frac18\le\frac{97}{1024};
\]
\eqref{eq:29} therefore gives net slope at least $223/1024$. The four margins in
\eqref{eq:margins} persist along the complete ranges $q\ge27$.

\section{Proof of Theorem~\ref{thm:main}}\label{sec:proof}

Fix $a\ge27$ or $a\le-27$ and write $\theta=\operatorname{Li}_3(1/a)$. Use
the corresponding seed and choose $p$ from \eqref{eq:19}. Take $m=12p^k$
on the positive side or $m=16p^k$ on the negative side. By
Lemma~\ref{lem:decomposition}, the determinant identity has the
form required by Fact~\ref{fact:criterion}. Define
\[
 \alpha(a)=A_+(\log|a|),\quad \tau(a)=\tau_+(|a|)\quad(a\ge27),
\]
and
\[
 \alpha(a)=A_-(\log|a|),\quad \tau(a)=\tau_-(|a|)\quad(a\le-27).
\]
Combine the coefficient integrality \eqref{eq:10}, the tower congruence
\eqref{eq:23}, the arithmetic limsup \eqref{eq:15} and its branch rates
\eqref{eq:17}. Add Propositions~\ref{prop:positive} and
\ref{prop:negative}, through their limsup statements \eqref{eq:24} and
\eqref{eq:28}. Together they give
\[
 \limsup_{k\to\infty}m^{-2}\log
 \left|v^mD_m(a)G_{m,a}(u/v)\right|
 \le -\alpha(a)+\tau(a)<0
\]
The endpoint bounds \eqref{eq:44} and the slope estimates immediately
following \eqref{eq:margins} give $\alpha(a)>\tau(a)$ throughout the two
ranges. The analytic limsups above hold for every $m=12Q$ and $m=16Q$,
respectively, hence also along $m=12p^k$ and $m=16p^k$. If
$\theta=u/v$ in lowest terms, the integer
$Z_m=v^mD_m(a)G_{m,a}(u/v)$ is nonzero by the tower congruence
\eqref{eq:23}. It is integral by \eqref{eq:10}. The factor $v^m$ is
only $\exp(O(m))$. Hence a tail of these nonzero integers has absolute
value strictly between zero and one. This is impossible, and the theorem
follows.



\appendix
\section{Positive density data}
\label{app:positive-data}
The positive endpoint certificates use the integer densities below. Cell $i$ is $[1/3+i/128,1/3+(i+1)/128)$. Each numerator is divided by $2^{40}$. The two columns belong to $q=27$ and $q=28$. Rows with $162\le i\le252$ are omitted; all their entries are zero.
\small
\renewcommand{\arraystretch}{1.05}
\begin{longtable}{@{}r r r@{\qquad\qquad}r r r@{}}
\caption{Positive density numerators; $n_{i,27}=n_{i,28}=0$ for $162\le i\le252$.}\label{tab:positive-density}\\
\toprule $i$ & $n_{i,27}$ & $n_{i,28}$ & $i$ & $n_{i,27}$ & $n_{i,28}$\\ \midrule
\endfirsthead
\toprule $i$ & $n_{i,27}$ & $n_{i,28}$ & $i$ & $n_{i,27}$ & $n_{i,28}$\\ \midrule
\endhead
0 & 3\,298\,534\,883\,328 & 3\,298\,534\,883\,328 & 83 & 622\,012\,586\,261 & 622\,050\,767\,912\\
1 & 3\,298\,534\,883\,328 & 3\,298\,534\,883\,328 & 84 & 614\,333\,078\,825 & 614\,287\,952\,151\\
2 & 3\,298\,534\,883\,328 & 3\,298\,534\,883\,328 & 85 & 606\,735\,467\,750 & 606\,608\,112\,256\\
3 & 3\,298\,534\,883\,328 & 3\,298\,534\,883\,328 & 86 & 599\,217\,211\,295 & 599\,008\,231\,565\\
4 & 3\,298\,534\,883\,328 & 3\,298\,534\,883\,328 & 87 & 591\,775\,970\,481 & 591\,484\,286\,035\\
5 & 3\,298\,534\,883\,328 & 3\,298\,534\,883\,328 & 88 & 584\,408\,472\,654 & 584\,032\,767\,761\\
6 & 3\,298\,534\,883\,328 & 3\,298\,534\,883\,328 & 89 & 577\,111\,889\,077 & 576\,651\,861\,460\\
7 & 2\,885\,832\,828\,696 & 2\,966\,263\,434\,188 & 90 & 569\,884\,096\,023 & 569\,339\,518\,896\\
8 & 2\,572\,916\,852\,719 & 2\,600\,553\,580\,905 & 91 & 562\,722\,486\,721 & 562\,093\,583\,640\\
9 & 2\,404\,674\,742\,681 & 2\,429\,959\,090\,809 & 92 & 555\,624\,724\,241 & 554\,911\,785\,762\\
10 & 2\,269\,717\,194\,177 & 2\,290\,223\,124\,257 & 93 & 548\,588\,660\,234 & 547\,791\,002\,500\\
11 & 2\,160\,719\,795\,372 & 2\,178\,694\,098\,863 & 94 & 541\,611\,628\,478 & 540\,729\,044\,235\\
12 & 2\,068\,303\,116\,386 & 2\,084\,358\,914\,107 & 95 & 534\,690\,937\,825 & 533\,724\,285\,633\\
13 & 1\,988\,106\,599\,617 & 2\,002\,705\,276\,844 & 96 & 527\,824\,999\,014 & 526\,773\,860\,583\\
14 & 1\,917\,217\,248\,650 & 1\,930\,647\,240\,393 & 97 & 521\,012\,316\,498 & 519\,874\,303\,491\\
15 & 1\,853\,688\,482\,146 & 1\,866\,157\,169\,433 & 98 & 514\,250\,142\,564 & 513\,023\,788\,088\\
16 & 1\,796\,134\,232\,221 & 1\,807\,790\,787\,793 & 99 & 507\,535\,377\,793 & 506\,221\,236\,102\\
17 & 1\,743\,532\,108\,475 & 1\,754\,487\,432\,635 & 100 & 500\,866\,279\,033 & 499\,464\,531\,627\\
18 & 1\,695\,103\,172\,824 & 1\,705\,444\,696\,477 & 101 & 494\,242\,009\,076 & 492\,751\,072\,173\\
19 & 1\,650\,239\,036\,412 & 1\,660\,039\,591\,631 & 102 & 487\,660\,717\,448 & 486\,079\,124\,444\\
20 & 1\,608\,458\,569\,067 & 1\,617\,775\,705\,817 & 103 & 481\,119\,346\,803 & 479\,446\,789\,252\\
21 & 1\,569\,373\,439\,666 & 1\,578\,250\,330\,129 & 104 & 474\,615\,344\,944 & 472\,851\,308\,457\\
22 & 1\,532\,662\,017\,100 & 1\,541\,136\,606\,734 & 105 & 468\,147\,362\,655 & 466\,290\,105\,005\\
23 & 1\,498\,056\,225\,875 & 1\,506\,164\,904\,576 & 106 & 461\,714\,229\,410 & 459\,761\,900\,645\\
24 & 1\,465\,333\,913\,300 & 1\,473\,106\,191\,037 & 107 & 455\,313\,811\,067 & 453\,265\,808\,862\\
25 & 1\,434\,307\,156\,504 & 1\,441\,764\,439\,830 & 108 & 448\,943\,409\,896 & 446\,799\,708\,779\\
26 & 1\,404\,811\,671\,588 & 1\,411\,975\,419\,573 & 109 & 442\,600\,868\,225 & 440\,360\,370\,311\\
27 & 1\,376\,705\,690\,614 & 1\,383\,597\,500\,366 & 110 & 436\,284\,765\,300 & 433\,944\,814\,059\\
28 & 1\,349\,869\,030\,884 & 1\,356\,504\,448\,938 & 111 & 429\,993\,546\,868 & 427\,551\,359\,822\\
29 & 1\,324\,196\,231\,791 & 1\,330\,587\,788\,739 & 112 & 423\,724\,966\,188 & 421\,179\,066\,889\\
30 & 1\,299\,591\,575\,199 & 1\,305\,753\,780\,226 & 113 & 417\,476\,547\,260 & 414\,826\,008\,656\\
31 & 1\,275\,971\,504\,575 & 1\,281\,916\,940\,006 & 114 & 411\,246\,311\,845 & 408\,489\,662\,946\\
32 & 1\,253\,264\,055\,032 & 1\,259\,001\,968\,334 & 115 & 405\,032\,602\,025 & 402\,168\,104\,928\\
33 & 1\,231\,402\,873\,061 & 1\,236\,943\,639\,366 & 116 & 398\,833\,314\,829 & 395\,859\,116\,566\\
34 & 1\,210\,327\,818\,537 & 1\,215\,680\,904\,025 & 117 & 392\,646\,208\,906 & 389\,560\,227\,707\\
35 & 1\,189\,986\,648\,507 & 1\,195\,158\,816\,234 & 118 & 386\,469\,149\,156 & 383\,268\,749\,568\\
36 & 1\,170\,331\,397\,942 & 1\,175\,329\,274\,750 & 119 & 380\,299\,767\,714 & 376\,981\,514\,354\\
37 & 1\,151\,317\,613\,315 & 1\,156\,148\,287\,118 & 120 & 374\,135\,763\,093 & 370\,696\,177\,386\\
38 & 1\,132\,905\,886\,597 & 1\,137\,575\,817\,319 & 121 & 367\,974\,952\,640 & 364\,410\,795\,140\\
39 & 1\,115\,060\,842\,055 & 1\,119\,574\,919\,777 & 122 & 361\,815\,086\,406 & 358\,124\,097\,999\\
40 & 1\,097\,748\,802\,721 & 1\,102\,112\,768\,050 & 123 & 355\,653\,493\,558 & 351\,833\,699\,890\\
41 & 1\,080\,938\,967\,357 & 1\,085\,157\,691\,959 & 124 & 349\,487\,525\,942 & 345\,535\,972\,733\\
42 & 1\,064\,604\,071\,733 & 1\,068\,681\,708\,434 & 125 & 343\,314\,177\,368 & 339\,227\,397\,144\\
43 & 1\,048\,717\,958\,913 & 1\,052\,659\,688\,523 & 126 & 337\,130\,791\,040 & 332\,903\,443\,659\\
44 & 1\,033\,257\,423\,905 & 1\,037\,066\,353\,123 & 127 & 330\,934\,714\,376 & 326\,561\,532\,878\\
45 & 1\,018\,200\,044\,902 & 1\,021\,880\,109\,083 & 128 & 324\,722\,903\,577 & 320\,198\,149\,644\\
46 & 1\,003\,525\,874\,709 & 1\,007\,078\,857\,138 & 129 & 318\,491\,787\,088 & 313\,810\,415\,083\\
47 & 989\,214\,576\,772 & 992\,644\,997\,276 & 130 & 312\,237\,660\,440 & 307\,394\,335\,365\\
48 & 975\,249\,426\,155 & 978\,560\,255\,853 & 131 & 305\,956\,916\,371 & 300\,946\,029\,728\\
49 & 961\,614\,613\,950 & 964\,808\,012\,985 & 132 & 299\,645\,605\,066 & 294\,461\,347\,829\\
50 & 948\,294\,135\,274 & 951\,372\,775\,929 & 133 & 293\,299\,436\,090 & 287\,934\,926\,304\\
51 & 935\,273\,085\,067 & 938\,239\,518\,829 & 134 & 286\,913\,758\,223 & 281\,360\,884\,753\\
52 & 922\,538\,698\,776 & 925\,394\,146\,872 & 135 & 280\,483\,988\,239 & 274\,732\,567\,154\\
53 & 910\,077\,910\,710 & 912\,824\,093\,089 & 136 & 274\,004\,651\,032 & 268\,044\,078\,975\\
54 & 897\,878\,310\,196 & 900\,519\,050\,516 & 137 & 267\,469\,791\,656 & 261\,289\,096\,273\\
55 & 885\,930\,295\,374 & 888\,467\,336\,853 & 138 & 260\,873\,025\,482 & 254\,460\,761\,612\\
56 & 874\,222\,852\,725 & 876\,656\,745\,923 & 139 & 254\,207\,328\,193 & 247\,551\,018\,043\\
57 & 862\,744\,015\,244 & 865\,078\,697\,765 & 140 & 247\,464\,701\,534 & 240\,550\,856\,530\\
58 & 851\,485\,807\,964 & 853\,723\,777\,364 & 141 & 240\,636\,654\,445 & 233\,449\,009\,424\\
59 & 840\,441\,243\,359 & 842\,580\,490\,827 & 142 & 233\,713\,639\,948 & 226\,232\,746\,699\\
60 & 829\,600\,075\,777 & 831\,640\,958\,545 & 143 & 226\,684\,558\,220 & 218\,888\,068\,958\\
61 & 818\,952\,623\,644 & 820\,899\,275\,764 & 144 & 219\,536\,950\,993 & 211\,399\,240\,675\\
62 & 808\,492\,968\,119 & 810\,345\,811\,167 & 145 & 212\,256\,823\,490 & 203\,748\,391\,803\\
63 & 798\,215\,511\,830 & 799\,972\,308\,091 & 146 & 204\,827\,860\,432 & 195\,914\,700\,575\\
64 & 788\,111\,907\,631 & 789\,775\,555\,841 & 147 & 197\,230\,645\,761 & 187\,873\,706\,515\\
65 & 778\,174\,308\,368 & 779\,749\,698\,740 & 148 & 189\,442\,528\,264 & 179\,595\,090\,436\\
66 & 768\,398\,074\,353 & 769\,884\,876\,401 & 149 & 181\,436\,911\,638 & 171\,041\,130\,629\\
67 & 758\,778\,213\,476 & 760\,174\,129\,545 & 150 & 173\,181\,443\,994 & 162\,165\,562\,957\\
68 & 749\,307\,875\,938 & 750\,614\,007\,649 & 151 & 164\,635\,705\,157 & 152\,909\,243\,355\\
69 & 739\,980\,958\,044 & 741\,200\,670\,769 & 152 & 155\,749\,296\,158 & 143\,194\,690\,447\\
70 & 730\,792\,985\,727 & 731\,927\,022\,986 & 153 & 146\,458\,187\,298 & 132\,918\,953\,924\\
71 & 721\,739\,398\,760 & 722\,787\,060\,370 & 154 & 136\,677\,474\,076 & 121\,939\,413\,447\\
72 & 712\,815\,494\,419 & 713\,776\,274\,750 & 155 & 126\,290\,632\,814 & 110\,040\,955\,919\\
73 & 704\,016\,555\,770 & 704\,891\,684\,344 & 156 & 115\,133\,084\,799 & 96\,902\,239\,046\\
74 & 695\,337\,727\,252 & 696\,128\,944\,904 & 157 & 102\,944\,846\,637 & 81\,791\,089\,438\\
75 & 686\,774\,873\,697 & 687\,482\,377\,327 & 158 & 89\,369\,632\,159 & 64\,286\,873\,668\\
76 & 678\,324\,907\,137 & 678\,946\,725\,136 & 159 & 73\,241\,560\,357 & 34\,652\,340\,814\\
77 & 669\,983\,602\,840 & 670\,519\,599\,327 & 160 & 55\,433\,749\,801 & 0\\
78 & 661\,746\,757\,665 & 662\,198\,323\,890 & 161 & 11\,861\,762\,313 & 0\\
79 & 653\,610\,372\,381 & 653\,980\,228\,320 & 253 & 10 & 6\\
80 & 645\,572\,367\,083 & 645\,860\,038\,638 & 254 & 1 & 1\\
81 & 637\,629\,226\,940 & 637\,834\,078\,443 & 255 & 1 & 0\\
82 & 629\,777\,087\,228 & 629\,898\,557\,487 &  &  & \\
\bottomrule
\end{longtable}
\normalsize
\noindent Each column has $0\le n_{i,q}\le3\cdot2^{40}$ and $\sum_i n_{i,q}=128\cdot2^{40}$. The entries are exact integer numerators.
\section{Negative rational density data}
\label{app:negative-data}
The negative endpoint certificates use the two rational densities below.
The parameter $u$ is arclength on the two branches
\[
 z_+(u)=\frac5{16}+\frac{4+i}{\sqrt{17}}u,\qquad
 z_-(u)=\frac5{16}+\frac{4-i}{\sqrt{17}}u,\qquad 0\le u\le3.
\]
For the cross-ray primitives put
$A=(4+i)/\sqrt{17}$ and $\beta=-\overline A=(-4+i)/\sqrt{17}$. Then
$Au+\beta v=z_+(u)-z_-(v)$. The second branch has direction $\overline A$.
The cell edges are
$u_j=3(j/256)^3$. Rows with $j\ge172$ are omitted; all their entries are zero.
The table entries are exact integers.
\small
\begin{longtable}{@{}r r r@{\qquad\qquad}r r r@{}}
\caption{Negative density counts, denominators $10^{12}-6$ and $10^{12}-8$; $c_{j,27}=c_{j,28}=0$ for $j\ge172$.}\label{tab:negative-density}\\
\toprule $j$ & $c_{j,27}$ & $c_{j,28}$ & $j$ & $c_{j,27}$ & $c_{j,28}$\\ \midrule
\endfirsthead
\toprule $j$ & $c_{j,27}$ & $c_{j,28}$ & $j$ & $c_{j,27}$ & $c_{j,28}$\\ \midrule
\endhead
0 & 140\,540\,525 & 142\,217\,998 & 86 & 7\,271\,281\,123 & 7\,303\,764\,615\\
1 & 261\,451\,321 & 264\,261\,111 & 87 & 7\,344\,376\,065 & 7\,376\,969\,350\\
2 & 363\,701\,611 & 367\,659\,490 & 88 & 7\,416\,401\,802 & 7\,449\,053\,094\\
3 & 453\,705\,003 & 458\,485\,706 & 89 & 7\,487\,335\,045 & 7\,519\,989\,402\\
4 & 536\,473\,939 & 542\,273\,883 & 90 & 7\,557\,144\,778 & 7\,589\,772\,163\\
5 & 614\,669\,706 & 621\,361\,877 & 91 & 7\,625\,796\,730 & 7\,658\,396\,089\\
6 & 689\,893\,588 & 697\,238\,114 & 92 & 7\,693\,233\,651 & 7\,725\,843\,275\\
7 & 763\,040\,641 & 770\,967\,704 & 93 & 7\,759\,430\,842 & 7\,792\,070\,509\\
8 & 834\,822\,469 & 843\,419\,489 & 94 & 7\,824\,335\,414 & 7\,857\,015\,320\\
9 & 905\,684\,947 & 914\,712\,009 & 95 & 7\,887\,930\,583 & 7\,920\,629\,973\\
10 & 976\,011\,761 & 985\,674\,694 & 96 & 7\,950\,207\,662 & 7\,982\,901\,909\\
11 & 1\,045\,990\,258 & 1\,056\,335\,435 & 97 & 8\,011\,161\,166 & 8\,043\,837\,986\\
12 & 1\,115\,838\,641 & 1\,126\,731\,236 & 98 & 8\,070\,799\,147 & 8\,103\,435\,277\\
13 & 1\,185\,744\,485 & 1\,197\,173\,719 & 99 & 8\,129\,080\,451 & 8\,161\,674\,235\\
14 & 1\,255\,906\,188 & 1\,267\,848\,817 & 100 & 8\,185\,957\,309 & 8\,218\,517\,762\\
15 & 1\,326\,378\,309 & 1\,338\,816\,011 & 101 & 8\,241\,427\,785 & 8\,273\,908\,536\\
16 & 1\,397\,237\,182 & 1\,410\,206\,063 & 102 & 8\,295\,441\,930 & 8\,327\,782\,491\\
17 & 1\,468\,565\,937 & 1\,482\,065\,631 & 103 & 8\,348\,017\,106 & 8\,380\,128\,377\\
18 & 1\,540\,470\,058 & 1\,554\,419\,086 & 104 & 8\,399\,056\,270 & 8\,431\,008\,634\\
19 & 1\,612\,985\,496 & 1\,627\,345\,506 & 105 & 8\,448\,533\,982 & 8\,480\,439\,151\\
20 & 1\,686\,109\,451 & 1\,700\,953\,931 & 106 & 8\,496\,439\,067 & 8\,528\,307\,528\\
21 & 1\,759\,886\,518 & 1\,775\,205\,628 & 107 & 8\,542\,791\,674 & 8\,574\,486\,950\\
22 & 1\,834\,332\,590 & 1\,850\,071\,630 & 108 & 8\,587\,618\,477 & 8\,618\,992\,632\\
23 & 1\,909\,514\,382 & 1\,925\,633\,468 & 109 & 8\,630\,829\,261 & 8\,661\,927\,493\\
24 & 1\,985\,414\,813 & 2\,001\,951\,334 & 110 & 8\,672\,314\,550 & 8\,703\,291\,099\\
25 & 2\,062\,008\,834 & 2\,078\,986\,141 & 111 & 8\,712\,110\,309 & 8\,742\,946\,397\\
26 & 2\,139\,339\,129 & 2\,156\,725\,739 & 112 & 8\,750\,176\,909 & 8\,780\,824\,766\\
27 & 2\,217\,420\,126 & 2\,235\,211\,994 & 113 & 8\,786\,530\,820 & 8\,816\,929\,623\\
28 & 2\,296\,229\,927 & 2\,314\,466\,234 & 114 & 8\,821\,148\,381 & 8\,851\,240\,883\\
29 & 2\,375\,762\,101 & 2\,394\,446\,069 & 115 & 8\,853\,965\,252 & 8\,883\,765\,766\\
30 & 2\,455\,985\,113 & 2\,475\,117\,286 & 116 & 8\,884\,986\,588 & 8\,914\,494\,208\\
31 & 2\,536\,916\,350 & 2\,556\,447\,255 & 117 & 8\,914\,160\,226 & 8\,943\,270\,999\\
32 & 2\,618\,559\,206 & 2\,638\,410\,477 & 118 & 8\,941\,455\,668 & 8\,969\,957\,130\\
33 & 2\,700\,907\,329 & 2\,721\,033\,529 & 119 & 8\,966\,786\,746 & 8\,994\,664\,508\\
34 & 2\,783\,913\,083 & 2\,804\,365\,072 & 120 & 8\,990\,150\,975 & 9\,017\,619\,226\\
35 & 2\,867\,560\,367 & 2\,888\,390\,055 & 121 & 9\,011\,513\,468 & 9\,038\,810\,379\\
36 & 2\,951\,843\,004 & 2\,973\,037\,533 & 122 & 9\,030\,797\,940 & 9\,057\,928\,634\\
37 & 3\,036\,727\,781 & 3\,058\,251\,294 & 123 & 9\,047\,959\,528 & 9\,074\,649\,583\\
38 & 3\,122\,175\,449 & 3\,144\,038\,906 & 124 & 9\,062\,954\,115 & 9\,088\,896\,328\\
39 & 3\,208\,172\,733 & 3\,230\,395\,933 & 125 & 9\,075\,777\,030 & 9\,100\,848\,019\\
40 & 3\,294\,699\,990 & 3\,317\,285\,043 & 126 & 9\,086\,389\,380 & 9\,110\,727\,562\\
41 & 3\,381\,718\,339 & 3\,404\,657\,728 & 127 & 9\,094\,662\,617 & 9\,118\,537\,364\\
42 & 3\,469\,217\,890 & 3\,492\,480\,725 & 128 & 9\,100\,537\,599 & 9\,123\,994\,734\\
43 & 3\,557\,179\,576 & 3\,580\,744\,413 & 129 & 9\,103\,976\,314 & 9\,126\,756\,271\\
44 & 3\,645\,583\,651 & 3\,669\,436\,031 & 130 & 9\,104\,931\,232 & 9\,126\,727\,580\\
45 & 3\,734\,354\,505 & 3\,758\,526\,323 & 131 & 9\,103\,261\,572 & 9\,124\,048\,078\\
46 & 3\,823\,492\,984 & 3\,847\,968\,982 & 132 & 9\,098\,842\,747 & 9\,118\,747\,158\\
47 & 3\,912\,928\,913 & 3\,937\,712\,428 & 133 & 9\,091\,576\,840 & 9\,110\,592\,597\\
48 & 4\,002\,643\,250 & 4\,027\,722\,788 & 134 & 9\,081\,425\,676 & 9\,099\,364\,692\\
49 & 4\,092\,610\,641 & 4\,117\,991\,994 & 135 & 9\,068\,283\,243 & 9\,085\,061\,206\\
50 & 4\,182\,796\,133 & 4\,208\,498\,764 & 136 & 9\,051\,951\,789 & 9\,067\,679\,190\\
51 & 4\,273\,196\,226 & 4\,299\,187\,192 & 137 & 9\,032\,309\,625 & 9\,046\,918\,113\\
52 & 4\,363\,741\,968 & 4\,389\,995\,946 & 138 & 9\,009\,221\,788 & 9\,022\,472\,399\\
53 & 4\,454\,405\,074 & 4\,480\,895\,428 & 139 & 8\,982\,617\,488 & 8\,994\,300\,776\\
54 & 4\,545\,159\,794 & 4\,571\,879\,737 & 140 & 8\,952\,238\,022 & 8\,962\,324\,092\\
55 & 4\,635\,953\,258 & 4\,662\,929\,300 & 141 & 8\,917\,780\,992 & 8\,926\,287\,605\\
56 & 4\,726\,742\,697 & 4\,753\,994\,578 & 142 & 8\,879\,031\,507 & 8\,885\,863\,474\\
57 & 4\,817\,490\,148 & 4\,845\,015\,105 & 143 & 8\,835\,875\,523 & 8\,840\,781\,200\\
58 & 4\,908\,147\,661 & 4\,935\,943\,324 & 144 & 8\,788\,123\,468 & 8\,790\,842\,965\\
59 & 4\,998\,696\,089 & 5\,026\,751\,567 & 145 & 8\,735\,314\,200 & 8\,735\,726\,253\\
60 & 5\,089\,119\,045 & 5\,117\,420\,789 & 146 & 8\,677\,085\,910 & 8\,675\,063\,681\\
61 & 5\,179\,384\,107 & 5\,207\,921\,943 & 147 & 8\,613\,120\,225 & 8\,608\,515\,497\\
62 & 5\,269\,426\,388 & 5\,298\,205\,557 & 148 & 8\,543\,105\,602 & 8\,535\,747\,731\\
63 & 5\,359\,217\,132 & 5\,388\,213\,135 & 149 & 8\,466\,593\,818 & 8\,456\,030\,313\\
64 & 5\,448\,713\,176 & 5\,477\,900\,666 & 150 & 8\,383\,030\,498 & 8\,368\,928\,198\\
65 & 5\,537\,874\,776 & 5\,567\,249\,859 & 151 & 8\,291\,808\,926 & 8\,274\,129\,435\\
66 & 5\,626\,668\,008 & 5\,656\,243\,489 & 152 & 8\,192\,433\,115 & 8\,170\,940\,771\\
67 & 5\,715\,065\,634 & 5\,744\,845\,063 & 153 & 8\,084\,283\,653 & 8\,058\,379\,664\\
68 & 5\,803\,012\,246 & 5\,833\,004\,641 & 154 & 7\,966\,570\,714 & 7\,935\,553\,073\\
69 & 5\,890\,480\,389 & 5\,920\,678\,097 & 155 & 7\,838\,343\,944 & 7\,801\,750\,381\\
70 & 5\,977\,447\,442 & 6\,007\,829\,650 & 156 & 7\,698\,603\,425 & 7\,656\,125\,755\\
71 & 6\,063\,860\,154 & 6\,094\,430\,692 & 157 & 7\,546\,277\,197 & 7\,497\,428\,193\\
72 & 6\,149\,679\,596 & 6\,180\,450\,129 & 158 & 7\,380\,133\,560 & 7\,324\,112\,186\\
73 & 6\,234\,900\,463 & 6\,265\,837\,571 & 159 & 7\,198\,670\,206 & 7\,134\,475\,246\\
74 & 6\,319\,460\,956 & 6\,350\,549\,141 & 160 & 7\,000\,072\,394 & 6\,926\,610\,257\\
75 & 6\,403\,343\,164 & 6\,434\,559\,295 & 161 & 6\,782\,275\,782 & 6\,698\,335\,400\\
76 & 6\,486\,508\,050 & 6\,517\,857\,236 & 162 & 6\,542\,881\,755 & 6\,446\,994\,517\\
77 & 6\,568\,907\,903 & 6\,600\,429\,001 & 163 & 6\,278\,760\,059 & 6\,169\,101\,464\\
78 & 6\,650\,533\,948 & 6\,682\,226\,511 & 164 & 5\,985\,945\,350 & 5\,860\,117\,045\\
79 & 6\,731\,345\,218 & 6\,763\,174\,621 & 165 & 5\,659\,351\,952 & 5\,514\,071\,290\\
80 & 6\,811\,305\,178 & 6\,843\,221\,978 & 166 & 5\,292\,033\,560 & 5\,122\,756\,282\\
81 & 6\,890\,362\,781 & 6\,922\,363\,261 & 167 & 4\,874\,213\,991 & 4\,673\,750\,846\\
82 & 6\,968\,507\,599 & 7\,000\,597\,226 & 168 & 4\,390\,896\,209 & 4\,149\,894\,406\\
83 & 7\,045\,692\,810 & 7\,077\,888\,090 & 169 & 3\,817\,831\,507 & 3\,504\,377\,317\\
84 & 7\,121\,915\,572 & 7\,154\,192\,159 & 170 & 3\,100\,186\,827 & 2\,746\,774\,945\\
85 & 7\,197\,117\,216 & 7\,229\,490\,450 & 171 & 2\,151\,224\,279 & 1\,052\,541\,758\\
\bottomrule
\end{longtable}
\normalsize
\noindent The interval boxes for the directed bounds are in the ancillary data
file \texttt{certificate\_data.json}.  The script
\texttt{check\_energy.py} recomputes the masses, caps, pair kernels, boxes,
and tails.
\section{Seed data}
\label{app:seed-data}
For the positive seed, $H_{211}$ has degree $211$ and its values modulo $29$ for residues $0,\ldots,28$ are
\[
 \begin{gathered}
(11, 15, 1, 13, 13, 23, 17, 9, 20, 10, 27, 19, 19, 12, 2,\\
 5, 9, 21, 16, 7, 16, 8, 2, 8, 19, 9, 17, 26, 6).
\end{gathered}
\]
For the negative seed, $H_{378}$ has degree $378$ and its values modulo $37$ for residues $0,\ldots,36$ are
\[
 \begin{gathered}
(28, 3, 19, 24, 31, 21, 36, 26, 27, 1, 30, 22, 21, 18, 11, 25, 31, 26, 36,\\
 11, 34, 19, 36, 4, 36, 24, 9, 8, 4, 33, 11, 30, 11, 34, 12, 24, 33).
\end{gathered}
\]
The complete ascending coefficient arrays are ancillary data. They are
in the files \texttt{H211\_exact\_seed.json} and
\texttt{H378\_exact\_seed.json}. Section~\ref{sec:nonvanishing} displays
the identities and the denominator factorizations.




\section{Computer checks}\label{app:checks}

Table~\ref{tab:checks-clean} lists the checks. Each script reads data
printed in this paper or in the ancillary tables. The scripts use the Python
standard library except where stated. The seed check
\texttt{check\_seed.py} requires \texttt{python-flint}; the energy and
contour interval checks require \texttt{mpmath}.
The script \texttt{check\_tau.py} does not certify every printed digit
independently. The finite check \texttt{check\_contour.py} diagnoses the
residue mechanism for an undifferentiated $F$. It also performs a separate
triple-zero check. It is not an $F''$ certificate and does not replace the
uniform contour proof.

\begin{table}[!ht]
\caption{Exact and directed-interval computations.}\label{tab:checks-clean}
\small
\renewcommand{\arraystretch}{1.15}
\begin{tabular}{@{}>{\footnotesize\ttfamily}l>{\raggedright\arraybackslash}p{0.57\textwidth}@{}}
\normalfont\small Script & Statement checked; result\\ \hline
check\_energy.py & \textbf{Directed interval replay:} recomputes the positive and negative certificate bounds from the printed density tables and interval data, including masses, caps, pair kernels, boxes, and tails. PASS.\\
check\_seed.py & \textbf{Exact algebra:} rebuilds both seed determinants, verifies the degree-bounded identities at 307 and 553 integer points, then checks the JSON denominator factorizations and displayed residue vectors. PASS.\\
check\_profile.py & \textbf{Exact algebra:} recomputes the two conservative baseline upper bounds, all ten low-layer increments for each ratio, and the resulting profile inequalities. PASS.\\
check\_tau.py & \textbf{Exact rational arithmetic:} checks the four endpoint margins, the two stated lower margin thresholds at $q=27$, and the identities $1/2-27/200=73/200$ and $5/16-97/1024=223/1024$ using rational logarithm enclosures; it does not certify the displayed energy decimals. PASS.\\
check\_contour.py & \textbf{Finite diagnostic:} residue-mechanism checks for an undifferentiated $F$, a separate triple-zero check, sine-bound and signed-contour consistency checks; it is not an $F''$ certificate and does not replace the uniform contour proof. PASS.\\
\end{tabular}
\end{table}

\section*{Statements and declarations}

\textbf{Competing interests.} The author has no competing interests.

\textbf{Funding.} No funding was received for this work.

\textbf{Data availability.} The scripts of Appendix~\ref{app:checks}, the
complete seed arrays, the interval boxes, and their compact outputs accompany
the paper as ancillary files. These files are part of the proof record.

\begin{thebibliography}{18}

\bibitem{Andreief}
C.~Andr\'eief, \emph{Note sur une relation entre les int\'egrales d\'efinies
des produits des fonctions}, M\'em. Soc. Sci. Phys. Nat. Bordeaux (3)
\textbf{2} (1886), 1--14.

\bibitem{Davenport}
H.~Davenport, \emph{Multiplicative Number Theory}, 3rd ed., revised by
H.~L.~Montgomery, Graduate Texts in Mathematics, vol.~74, Springer, 2000.

\bibitem{Fauzan}
A.~Fauzan, \emph{$\zeta(5)$ is irrational}, Zenodo (2026),
\url{https://doi.org/10.5281/zenodo.22826419}.

\bibitem{Hata}
M.~Hata, \emph{On the linear independence of the values of polylogarithmic
functions}, J. Math. Pures Appl. (9) \textbf{69} (1990), 133--173.

\bibitem{Krattenthaler}
C.~Krattenthaler, \emph{Advanced determinant calculus}, S\'em. Lothar.
Combin. \textbf{42} (1999), Art.~B42q, 67~pp.

\bibitem{Lai}
L.~Lai, \emph{The linear independence of $1,\zeta(3),\zeta(4)$},
arXiv:2610.06838 (2026).

\bibitem{Liu}
S.-Q.~Liu, \emph{Li2Rational}, GitHub repository (2026), commit
\texttt{b2a61c4}, \url{https://github.com/siqiliu-tsinghua/Li2Rational}.

\bibitem{Marcovecchio}
R.~Marcovecchio, \emph{Linear independence of polylogarithms at algebraic
points}, Mosc. J. Comb. Number Theory \textbf{6} (2016), no.~2--3, 208--232.

\bibitem{Panzone}
P.~A.~Panzone, \emph{A simple proof of the irrationality of the trilog},
Rev. Un. Mat. Argentina \textbf{42} (2001), no.~2, 103--111.

\bibitem{Prellberg}
T.~Prellberg, \emph{Three irrationality results for the dilogarithm},\\
arXiv:\allowbreak 2609.06895 (2026).

\bibitem{SaffTotik}
 E.~B.~Saff and V.~Totik, \emph{Logarithmic Potentials with External
 Fields},
Grundlehren der mathematischen Wissenschaften, vol.~316, Springer, 1997.

\bibitem{TanZhang}
H.~Tan and Y.~Zhang, \emph{Uniform cubic bounds for irrational values of the
dilogarithm}, Zenodo (2026),
\url{https://doi.org/10.5281/zenodo.23247332}.

\bibitem{Zh}
R.~Zhang, \emph{Irrationality of the dilogarithm at reciprocals of
integers}, preprint (2026), Zenodo,
\url{https://doi.org/10.5281/zenodo.23296412}.

\bibitem{ZhB}
R.~Zhang, \emph{Irrationality of the trilogarithm at reciprocals of large
integers}, preprint (2026), Zenodo,
\url{https://doi.org/10.5281/zenodo.23329357}.

\end{thebibliography}

\end{document}

